% The Axon User Manual. Build with tools\build_manual.cmd, which first runs
% the scripts in examples\ (they write the figures into figures\), then
% latexmk. Every Axon listing in this file is a complete program that runs.
\documentclass[11pt,a4paper,openany]{report}

\usepackage[margin=2.5cm]{geometry}
\usepackage[T1]{fontenc}
\usepackage[utf8]{inputenc}
\usepackage{lmodern}
\usepackage[scaled=0.88]{beramono}
\usepackage{xcolor}
\usepackage{listings}
\usepackage{graphicx}
\usepackage{booktabs}
\usepackage{longtable}
\usepackage{tabularx}
\usepackage{array}
\usepackage[shortlabels]{enumitem}
\usepackage{caption}
\usepackage[colorlinks=true,linkcolor=blue!50!black,urlcolor=blue!50!black,
            citecolor=blue!50!black]{hyperref}

\setlist{itemsep=2pt,topsep=4pt}
\setlength{\parskip}{4pt}
\setlength{\parindent}{0pt}
\setlength{\emergencystretch}{3em}
\captionsetup{font=small,labelfont=bf}

% Colours of the Axon IDE's light theme.
\definecolor{axkeyword}{HTML}{0033B3}
\definecolor{axbuiltin}{HTML}{00796B}
\definecolor{axstring}{HTML}{A31515}
\definecolor{axnumber}{HTML}{098658}
\definecolor{axcomment}{HTML}{6A737D}
\definecolor{axdoc}{HTML}{8A5A00}
\definecolor{axback}{HTML}{F6F8FA}
\definecolor{axframe}{HTML}{D0D7DE}
\definecolor{outback}{HTML}{FFFFFF}

\lstdefinelanguage{Axon}{
  morekeywords={function,end,return,if,elseif,else,for,in,while,break,
    continue,local,global,try,catch,import,from,as,export,extern,record,
    true,false,nothing,and,or,not},
  morekeywords=[2]{zeros,ones,eye,rand,linspace,size,sum,prod,mean,all,any,
    dot,adjoint,isequal,keep,eltype,reshape,sparse,speye,nnz,full,splu,
    append,get,keys,haskey,dict,print,println,error,string,length,typeof,
    sqrt,exp,log,sin,cos,tan,asin,acos,atan,sinh,cosh,tanh,asinh,acosh,atanh,
    abs,int,floor,ceil,round,real,imag,conj,arg,complex,div,mod,min,max,
    factorial,format,pi,sort,sortperm,atexit},
  sensitive=true,
  morecomment=[l]{\#},
  morestring=[b]",
}

\lstdefinestyle{axonstyle}{
  language=Axon,
  basicstyle=\ttfamily\small,
  keywordstyle=\color{axkeyword}\bfseries,
  keywordstyle=[2]\color{axbuiltin},
  stringstyle=\color{axstring},
  commentstyle=\color{axcomment}\itshape,
  backgroundcolor=\color{axback},
  frame=single, rulecolor=\color{axframe}, framesep=4pt,
  xleftmargin=6pt, xrightmargin=6pt,
  columns=fullflexible, keepspaces=true,
  showstringspaces=false, upquote=true,
  aboveskip=8pt, belowskip=6pt,
  breaklines=true, breakatwhitespace=true,
  tabsize=2,
  literate={*}{{\char42}}1 {-}{{\char45}}1,
}
\lstdefinestyle{outstyle}{
  language={}, keywordstyle={}, stringstyle={}, commentstyle={},
  literate={*}{{\char42}}1 {-}{{\char45}}1,
  basicstyle=\ttfamily\small\color{black!75},
  backgroundcolor=\color{outback},
  frame=leftline, rulecolor=\color{axframe}, framerule=2pt,
  xleftmargin=10pt, framesep=6pt,
  columns=fullflexible, keepspaces=true, upquote=true,
  aboveskip=0pt, belowskip=8pt,
  escapeinside={(*@}{@*)},
}
\lstdefinestyle{shellstyle}{
  language={}, keywordstyle={}, stringstyle={}, commentstyle={},
  basicstyle=\ttfamily\small,
  backgroundcolor=\color{axback},
  frame=single, rulecolor=\color{axframe}, framesep=4pt,
  xleftmargin=6pt, xrightmargin=6pt,
  columns=fullflexible, keepspaces=true, upquote=true,
}
% axon: a complete program (checked by running it). axonout: what it prints.
% axonerr: a program that stops with an error; the following axonout is the
% error report. shell: a command line.
\lstnewenvironment{axon}{\lstset{style=axonstyle}}{}
\lstnewenvironment{axonerr}{\lstset{style=axonstyle}}{}
\lstnewenvironment{axonout}{\lstset{style=outstyle}}{}
\lstnewenvironment{shell}{\lstset{style=shellstyle}}{}
% axonrepl: a REPL session, checked with axon --transcript.
\lstnewenvironment{axonrepl}{\lstset{style=outstyle,aboveskip=8pt}}{}
\lstset{style=axonstyle}
\newcommand{\ex}[1]{\lstinputlisting[style=axonstyle,
  title={\small\ttfamily\detokenize{#1}}]{examples/#1}}
% Inline code, read verbatim: \code{x + 1}, or \code|{a, b}| with braces inside.
\newcommand{\code}{\lstinline[basicstyle=\ttfamily,keywordstyle=,
  keywordstyle={[2]},stringstyle=,commentstyle=,
  literate={*}{{\char42}}1 {-}{{\char45}}1]}
\newcommand{\menu}[1]{\textsf{#1}}
\newcommand{\key}[1]{\fbox{\footnotesize\sffamily #1}}
\newcommand{\fig}[3]{\begin{figure}[htbp]\centering
  \includegraphics[width=#2\linewidth]{figures/#1}\caption{#3}\label{fig:#1}
  \end{figure}}

\title{\Huge\bfseries Axon\\[6pt]\LARGE User Manual\\[18pt]
  \large The language and the IDE}
\author{Herbert Sauro\\[4pt]\small developed in collaboration with Claude Code}
\date{Version 0.1, November 2026}

\begin{document}
\pagenumbering{roman}
\maketitle
\tableofcontents
\clearpage
\pagenumbering{arabic}

% ===================================================================
\chapter{Introduction}

Axon is a programming language for scientific computing. Vectors, matrices,
complex numbers, differential equation solvers and plots are part of the
language and its standard library rather than add-ons, and the syntax borrows
what scientists already read fluently: MATLAB-style matrix literals, 1-based
indexing and \code{end}-delimited blocks. To that it adds what larger programs
need: modules with explicit exports, records, keyword arguments and optional
type annotations.

Axon comes with an IDE: an editor with syntax colouring and checks as you type,
a REPL that keeps its variables, help on any name with \key{F1}, a searchable
documentation viewer, and a plot pane in which line, bar, histogram, surface and
3D plots can be zoomed, panned, rotated and edited.

A short program gives the flavour:

\begin{axon}
import ode

# Linear algebra is built in: \ solves A x = b.
A = [4.0, 1.0; 1.0, 3.0]
b = [1.0; 2.0]
println(A \ b)

# Exponential decay, y' = -0.5 y, solved numerically.
decay = function(t, y) return -0.5 * y end
sol = ode.solve(decay, 2.0, [0, 10], rtol = 1e-8)
println("y(3) =", format(sol.at(3.0)[1], 6), " exact:", format(2 * exp(-1.5), 6))
\end{axon}
\begin{axonout}
2-vector{real}
  0.0909091
   0.636364
y(3) = 0.446260  exact: 0.446260
\end{axonout}

\section*{About this manual}

This manual teaches Axon and its IDE to someone who knows another numerical
language, such as MATLAB, Python with NumPy, R or Julia. Chapter~\ref{ch:start}
gets a first program running; Chapter~\ref{ch:ide} tours the IDE;
Chapters~\ref{ch:lang} to \ref{ch:c} cover the language, the numerical modules,
working with text, files and data, plotting, writing modules of your own, performance and calling
C. Appendix~\ref{app:builtins} lists every
built-in function, and Appendix~\ref{app:examples} prints the example programs
that come with Axon.

Every program in a shaded box is complete and runs as printed; what it prints
follows it, with a bar on the left. The longer examples are files in the
manual's \texttt{examples} folder, and the figures in Chapter~\ref{ch:plot} are
drawn by those files when the manual is built.

The \emph{Axon Language Reference}, in the IDE under \menu{Help > Language
Reference}, is the authority on every rule; this manual explains and
illustrates, and points to the reference where the details matter.

\section*{Acknowledgement}

Axon, its IDE and this manual were developed by Herbert Sauro in collaboration
with Claude Code, Anthropic's AI coding assistant.

% ===================================================================
\chapter{Getting started}\label{ch:start}

\section{What you get}

Axon consists of two programs and some folders that live together:

\begin{itemize}
\item \texttt{axon} (\texttt{axon.exe} on Windows), the interpreter. It runs
  programs, starts the REPL and does everything else from a terminal.
\item \texttt{AxonIDE}, the IDE, which runs \texttt{axon} for you.
\item \texttt{stdlib}, the standard library (\texttt{linalg}, \texttt{ode},
  \texttt{plot} and others, written in Axon); \texttt{examples}, the example
  programs; and \texttt{help}, the documentation the IDE shows.
\item The native libraries: OpenBLAS for linear algebra, SUNDIALS CVODE for
  stiff differential equations, Skia for graphics and MIR for the JIT compiler.
\end{itemize}

Axon runs on Windows (x64) and macOS (Apple Silicon). On macOS everything is
inside the \texttt{AxonIDE.app} bundle.

\paragraph{Building from source.} Axon is written in Delphi and built with
Embarcadero RAD Studio 13. Open \texttt{Axon.groupproj} and use
\menu{Build All}, or from a command prompt:

\begin{shell}
build.cmd Release
\end{shell}

The interpreter and the IDE land in \texttt{Win64\textbackslash Release},
together with the libraries, \texttt{stdlib}, the examples and the help. The
IDE also needs four editor and chart components from the author's
repositories; the \texttt{README} in the source tree lists them.

\section{Your first program}

In the IDE, choose \menu{File > New} (\key{Ctrl+N}), type a program and press
\key{F9} to run it. The IDE asks you to save it first; Axon files end in
\texttt{.sci}.

\begin{axon}
println("Hello, world")
A = [1, 2; 3, 4]
println(A * A)
\end{axon}
\begin{axonout}
Hello, world
2(*@\texttimes@*)2 matrix{int}
   7  10
  15  22
\end{axonout}

The output appears in the \textbf{Output} pane. From a terminal, the same file
runs with

\begin{shell}
axon hello.sci
\end{shell}

A program that stops with an error prints the kind of error, a message and a
\emph{trace}: where it happened, innermost first. In the IDE the line is
underlined in the editor and each trace line is a link.

\begin{axonerr}
function lookup(l, k)
  return l[k]
end
names = {"S1", "S2", "S3"}
println(lookup(names, 5))
\end{axonerr}
\begin{axonout}
IndexError: index 5 is out of bounds for a dimension of length 3
  at program.sci:2:12 in lookup
  at program.sci:5:9
\end{axonout}

\section{Trying things out: the REPL}

The REPL runs one statement at a time and keeps your variables between them.
In the IDE it is the \textbf{REPL} tab at the bottom; from a terminal, run
\texttt{axon} with no file. Type an expression and press \key{Enter}; its value
is shown at once. A statement that continues over several lines, such as a
function, shows a \texttt{...} prompt until it is complete.

\begin{axonrepl}
> x = linspace(0.0, 1.0, 3)
x = 3-vector{real}
  0.0
  0.5
  1.0
> sum(x .^ 2)
1.25
> s = "Axon";
> length(s)
4
\end{axonrepl}

A trailing \code{;} suppresses the display. Two lines are commands of the REPL
rather than statements:

\begin{itemize}
\item \code{help name} shows the parameters and description of a function,
  record or module, for instance \code{help sum} or \code{help linalg.solve}.
  \code{help} alone lists the commands.
\item \code{run "file.sci"} runs a whole file in the session, so that its
  variables stay afterwards for you to inspect.
\end{itemize}

\section{Finding help}

\begin{itemize}
\item In the IDE, put the caret on a name, such as \code{sin} or
  \code{ode.solve}, and press \key{F1}. The \textbf{Help} pane shows the
  parameters and description, followed by the section of the Language Reference
  that documents it.
\item \menu{Help > Documentation} opens the IDE's guides, \menu{Help > Language
  Reference} the complete description of the language, and \menu{Help >
  Built-in Functions} the list of every function you can use without an
  \code{import}.
\item The examples on the IDE's \textbf{Welcome} tab are complete programs to
  open, change and run.
\end{itemize}

Functions you write yourself get help too: a comment that starts with \code{##}
directly above a function is its documentation, shown by \key{F1} and by
\code{help} at the REPL.

\begin{axon}
## The area of a circle of radius r.
function area(r)
  return pi * r^2
end
println(area(2.0))
\end{axon}
\begin{axonout}
12.5663706143592
\end{axonout}

% ===================================================================
\chapter{The IDE}\label{ch:ide}

\section{The window}

The IDE's window has four areas:

\begin{center}
\begin{tabularx}{\linewidth}{@{}lX@{}}
\toprule
Area & What it shows \\
\midrule
Editor tabs & Your files, with syntax colouring and line numbers. The
  \textbf{Welcome} tab lists recent files and the examples. \\
Plot pane (right) & The figures your programs show, with a history to step
  back through. \\
Output (bottom) & What a program prints, and its errors in red. \\
REPL (bottom) & An interactive session that keeps its variables. \\
Help (bottom) & Help on the word under the caret (\key{F1}). \\
\bottomrule
\end{tabularx}
\end{center}

The panes can be resized by dragging their edges. The IDE remembers the
window's size and position (including full screen on macOS), the pane sizes and
the recent files between sessions.

\paragraph{Tabs.} Close a tab with the \texttimes{} at its right end, a
middle-click, or \key{Ctrl+W}. The examples on the Welcome tab open as your own
copies, in a folder \texttt{Axon Examples} in your documents, so the originals
stay unchanged.

\paragraph{Copying output.} To copy from the Output or REPL pane, click a line
(Shift- or Ctrl-click for more) and press \key{Ctrl+C}, or right-click for
\menu{Copy} and \menu{Copy All}. With no line selected, Copy takes the whole
pane.

\section{Writing and running programs}

\key{F9} (\menu{Run > Run}) saves the current file and runs it with
\texttt{axon}, in the file's own folder, so \code{import} finds modules next to
it and files the program writes land there. \textbf{Stop} (\key{Ctrl+F2}) ends a
program that runs too long.

While you type, the IDE checks the syntax: about half a second after you stop, a
syntax error gets a wavy red underline, and hovering over it shows the message.
When a program stops with an error, the error's line is underlined and the caret
moves there; in the Output pane each trace line is a link to its place.

The grey lines in the Output pane are the IDE's own notes: \emph{Running \ldots}
and \emph{Finished in \ldots} around each run, and the JIT report after it
(Section~\ref{sec:ide-jit}). \menu{Edit > Preferences} can hide them.

The editor has the usual commands under \menu{Edit}: undo and redo, find and
replace (\key{Ctrl+F}), go to line (\key{Ctrl+G}), and \menu{Toggle Comment}
(\key{Ctrl+/}), which comments or uncomments the selected lines.

\section{The REPL pane}

Type a statement and press \key{Enter}. A session starts the first time you use
it, in the current file's folder, and keeps its variables until you press
\textbf{Restart}, which also stops a statement that does not finish.
\textbf{Clear} empties the transcript, and \key{Up} and \key{Down} recall
earlier lines.

\key{F9} runs a program on its own, and its variables are gone when it ends. To
keep them, use the \textbf{\textrightarrow\ REPL} button next to Run
(\menu{Run > Run in REPL}, \key{Shift+F9}): it saves the file and runs it in the
REPL session, so that afterwards you can inspect any of its variables there,
such as \code{sol.t} or \code{A}. \key{Ctrl+Enter} sends the current line, or
the selection, to the REPL, which is a quick way to try part of a program.

\section{Help and documentation}

\key{F1} shows help on the name under the caret: a built-in function, a
function from an imported module (\code{ode.solve}, or \code{solve} after
\code{from linalg import solve}), or a function or record in the file with a
\code{##} comment above it. Below the short description, the Help pane shows
the section of the Language Reference that documents the name; its last link
opens that section in the \textbf{Documentation} tab.

The Documentation tab (\menu{Help > Documentation}, \key{Ctrl+F1}) shows the
guides, the Language Reference and the white paper, with links between them,
\textbf{Back} and \textbf{Home} buttons, and syntax-coloured code.
\key{Ctrl+F} there opens a find bar: type to highlight every match, and use
\key{Enter} and \key{Shift+Enter} for the next and previous one.

\section{Preferences and appearance}

\menu{Edit > Preferences} (\key{Ctrl+,}) sets:

\begin{itemize}
\item the editor's font size;
\item the font size of the Output, REPL and Help panes;
\item the font size of the Documentation tab;
\item whether a run writes its notes (\emph{Running \ldots}, \emph{Finished in
  \ldots}) and the JIT report to the Output pane.
\end{itemize}

\menu{View > Dark Theme} switches the whole IDE, editor and documentation
included, to a dark theme, and back. The choices are kept between sessions.

\section{The plot pane}\label{sec:plotpane}

\code{plot.show()} displays the current figure in the plot pane, even while
the program is still running, and a program that ends without calling it shows
its figure anyway (Chapter~\ref{ch:plot}). The pane keeps every figure shown
since the IDE started, the newest 50:

\begin{itemize}
\item \textbf{\textleftarrow} and \textbf{\textrightarrow} in its header step
  through the figures, and the count between them shows which one you are
  looking at.
\item \textbf{Copy} puts the figure on the clipboard as an image, and
  \textbf{Save\ldots} writes it as PNG, PDF or SVG (PDF and SVG as vector
  graphics). Both take the figure as it looks now, with your zooming, panning,
  rotating and edits, and both take the whole figure, every cell of a grid.
\item \textbf{Clear} forgets them all.
\end{itemize}

\paragraph{Interactive charts.} Line, scatter, bar and histogram plots are drawn
in an interactive chart:

\begin{center}
\begin{tabular}{@{}ll@{}}
\toprule
Action & Effect \\
\midrule
Mouse wheel & Zoom in or out around the pointer \\
Drag & Pan \\
\key{Shift}+drag & Zoom to the box you draw \\
\key{Ctrl}+double-click & Back to the whole plot \\
Double-click & Edit titles, axes, series styles and the legend \\
Click a legend entry & Hide or show that series \\
Right-click & Copy the image, or export it as PNG or PDF \\
\bottomrule
\end{tabular}
\end{center}

3D plots (surfaces, \code{line3} and \code{scatter3}) are drawn in an interactive
3D chart:

\begin{center}
\begin{tabular}{@{}ll@{}}
\toprule
Action & Effect \\
\midrule
Drag & Rotate \\
Right-drag, or \key{Shift}+drag & Pan \\
Mouse wheel & Zoom \\
Hover & Show the coordinates of the nearest point \\
Double-click & Edit titles, axes, series styles, lighting and the legend \\
\key{Ctrl}+double-click & Back to the starting view \\
Arrow keys, \key{+}, \key{-}, \key{R}, \key{1}--\key{4} & Rotate, zoom, reset,
  preset views \\
Right-click & Copy the image, export PNG or PDF, and more \\
\bottomrule
\end{tabular}
\end{center}

A figure divided with \code{plot.subplot} is shown as a grid of such charts, one
per cell, each of which can be zoomed, panned or rotated and edited on its own.
Widen the plot pane by dragging its left edge to give the cells more room. A
chart's own right-click menu copies or exports just that cell; \textbf{Copy}
and \textbf{Save\ldots} in the pane's header take the whole grid. Changes made
in a chart are not kept when you step to another figure.

\section{Compiled and interpreted}\label{sec:ide-jit}

Functions are compiled to native code at their first call
(Chapter~\ref{ch:perf}). The \textbf{JIT} button in the toolbar, or
\menu{Run > Use JIT Compiler}, switches this off and on, so you can run the
same program both ways and compare: the results are the same, only the time
differs. After a run with the JIT on, the Output pane lists the functions that
were compiled and, for those that ran interpreted, the reason. The REPL follows
the setting when it starts; use \menu{Run > Restart REPL} after switching.

\section{Keyboard shortcuts}

\begin{center}
\begin{tabular}{@{}ll@{}}
\toprule
Key & Action \\
\midrule
\key{F9} & Run the current file \\
\key{Shift+F9} & Run the current file in the REPL, keeping its variables \\
\key{Ctrl+F2} & Stop the running program \\
\key{Ctrl+Enter} & Send the current line or selection to the REPL \\
\key{F1} & Help on the word under the caret \\
\key{Ctrl+F1} & The Documentation tab \\
\key{Ctrl+N}, \key{Ctrl+O}, \key{Ctrl+S} & New, open, save \\
\key{Shift+Ctrl+S} & Save as \\
\key{Ctrl+W} & Close the current tab \\
\key{Ctrl+F} & Find and replace; in the Documentation tab, find in the page \\
\key{Ctrl+G} & Go to line \\
\key{Ctrl+/} & Comment or uncomment the selected lines \\
\key{Ctrl+Z}, \key{Ctrl+Y} & Undo, redo \\
\key{Ctrl+,} & Preferences \\
\bottomrule
\end{tabular}
\end{center}

On macOS, use \key{Cmd} in place of \key{Ctrl}.

% ===================================================================
\chapter{The language}\label{ch:lang}

This chapter goes through the language from values to modules. Each section
explains the common case with examples and mentions the rules that most often
matter; the Language Reference has the rest.

\section{Programs, statements and comments}

A program is a sequence of statements, one per line. A \code{;} also ends a
statement, so several can share a line. An expression continues onto the next
line after a binary operator, a comma or the \code{=} of an assignment, and
inside brackets line ends are ignored altogether.

\begin{axon}
# A comment runs to the end of the line.
x = 1; y = 2          # two statements on one line
total = x +
        y             # continues after the operator
#= A block comment
   can span lines, and #= nest =#. =#
println(total)
\end{axon}
\begin{axonout}
3
\end{axonout}

Blocks (\code{if}, \code{for}, \code{while}, \code{try}, \code{function},
\code{record}) all close with \code{end}. Names are case-sensitive, and the
words of Table~\ref{tab:keywords} are reserved.

\begin{table}[htbp]
\centering
\caption{The reserved words.}\label{tab:keywords}
\begin{tabular}{@{}*{7}{l}@{}}
\toprule
\texttt{function} & \texttt{end} & \texttt{return} & \texttt{if} & \texttt{elseif}
  & \texttt{else} & \texttt{for} \\
\texttt{in} & \texttt{while} & \texttt{break} & \texttt{continue} & \texttt{local}
  & \texttt{global} & \texttt{try} \\
\texttt{catch} & \texttt{import} & \texttt{from} & \texttt{as} & \texttt{export}
  & \texttt{extern} & \texttt{record} \\
\texttt{true} & \texttt{false} & \texttt{nothing} & \texttt{and} & \texttt{or}
  & \texttt{not} & \\
\bottomrule
\end{tabular}
\end{table}

\section{Values and types}

Every value has a type, which \code{typeof} names:

\begin{center}
\begin{tabular}{@{}lll@{}}
\toprule
Type & Examples & Notes \\
\midrule
\texttt{int} & \code|42|, \code|-7|, \code|1_000_000| & 64-bit integers \\
\texttt{real} & \code|2.5|, \code|1.0e-3|, \code|6.02e23| & double precision \\
\texttt{complex} & \code|3 + 4i|, \code|2.5i| & \\
\texttt{bool} & \code|true|, \code|false| & \\
\texttt{string} & \code|"text"| & \\
\texttt{nothing} & \code|nothing| & no value \\
\texttt{vector}, \texttt{matrix} & \code|[1; 2; 3]|, \code|[1, 2; 3, 4]| &
  dense arrays \\
\texttt{range} & \code|1:10|, \code|0:0.1:1| & lazy sequences \\
\texttt{list} & \code|{1, "two", 3.0}| & any mix of values \\
\texttt{record} & \code|{a: 1.0, b: "x"}| & named fields \\
\texttt{dict} & \code|dict("Km" => 0.5)| & keys to values \\
\texttt{sparse} & \code|sparse(i, j, v, m, n)| & sparse matrices \\
\texttt{function} & \code|function(x) return 2*x end| & functions are values \\
\bottomrule
\end{tabular}
\end{center}

\begin{axon}
println(typeof(3), typeof(3.0), typeof(3i), typeof("3"), typeof([3; 4]),
        typeof({3}), typeof({x: 3}), typeof(1:3), typeof(nothing))
\end{axon}
\begin{axonout}
int real complex string vector list record range nothing
\end{axonout}

A real literal needs a digit on both sides of the point: \code{0.5} and
\code{1.0}, not \code{.5} or \code{1.}. That keeps \code{A.*B} and \code{1.+x}
unambiguous.

\paragraph{Value semantics.} Every value behaves as if it were copied when it is
assigned or passed to a function: after \code{B = A}, changing \code{B} never
changes \code{A}. The copy is made lazily, only when one of the two is changed,
so passing a large matrix to a function costs nothing.

\begin{axon}
A = [1.0, 2.0; 3.0, 4.0]
B = A
B[1, 1] = 100.0
println(A[1, 1], B[1, 1])
\end{axon}
\begin{axonout}
1.0 100.0
\end{axonout}

\paragraph{Type annotations.} A variable, a parameter or a function's result may
carry a type, which is then checked. A value of a narrower numeric type is
widened to fit, so \code{x: real = 0} stores \code{0.0}; a real where an
\code{int} is declared is an error.

\begin{axonerr}
n: int = 10
x: real = 0        # widened to 0.0
println(n, x)
n = 2.5            # a real is not an int
\end{axonerr}
\begin{axonout}
10 0.0
TypeError: variable "n" must be int, got real
  at program.sci:4:1
\end{axonout}

\section{Numbers}

Integers and reals are distinct. Arithmetic between them gives a real, and
\code{/} always does; \code{div} and \code{mod} are integer division and
remainder. Integer overflow is an error, never a silent wrap-around.

\begin{axon}
a = 7
b = 2
println(a / b, div(a, b), mod(a, b), div(-7, 2), mod(-7, 2))
println(2^10, 2^-1, 1.5e3, 0.1 + 0.2)
println(int(2.7), round(2.5), floor(-2.5), ceil(2.1))
\end{axon}
\begin{axonout}
3.5 3 1 -4 1
1024 0.5 1500.0 0.3
2 3 -3 3
\end{axonout}

Arithmetic promotes along \texttt{bool} \textrightarrow{} \texttt{int}
\textrightarrow{} \texttt{real} \textrightarrow{} \texttt{complex} and never the
other way. Real arithmetic follows IEEE 754: \code{1.0/0.0} is \code{Inf} and
\code{0.0/0.0} is NaN, and real functions stay real outside their domain, so
\code{sqrt(-4.0)} is NaN while \code{sqrt(-4 + 0i)} is \code{2i}.

\begin{axon}
z = 3 + 4i
println(abs(z), conj(z), real(z), imag(z), arg(1i))
println(sqrt(-4.0), sqrt(-4 + 0i), 1.0 / 0.0)
println(pi, e, exp(1.0) == e)
\end{axon}
\begin{axonout}
5.0 3.0 - 4.0i 3.0 4.0 1.5707963267949
NaN 0.0 + 2.0i Inf
3.14159265358979 2.71828182845905 true
\end{axonout}

Reals display with up to 15 significant digits, which is why \code{0.1 + 0.2}
shows as \code{0.3}; Section~\ref{sec:printing} shows how to choose the digits.

\section{Vectors and matrices}\label{sec:arrays}

Square brackets build arrays. Entries in a row are separated by commas, rows by
semicolons. A literal with one entry per row is a \emph{vector}, a column;
anything else is a \emph{matrix}. So \code{[1; 2; 3]} is a vector of length 3
and \code{[1, 2, 3]} is a $1\times3$ matrix, a row.

\begin{axon}
v = [1.0; 2.0; 3.0]
A = [2.0, 1.0, 0.0;
     1.0, 3.0, 1.0;
     0.0, 1.0, 4.0]
println(A * v)          # matrix times vector
println(v' * v)         # ' transposes: a row times a column is a scalar
println(v .* v)         # elementwise product
\end{axon}
\begin{axonout}
3-vector{real}
   4.0
  10.0
  14.0
14.0
3-vector{real}
  1.0
  4.0
  9.0
\end{axonout}

The plain operators \code{*}, \code{/}, \code{\} and \code{^} are the matrix
operations; the dotted ones, \code{.*}, \code{./}, \code{.\} and \code{.^}, work
element by element. \code{+} and \code{-} are always elementwise.

\paragraph{Making arrays.} \code{zeros(m, n)}, \code{ones(m, n)} and
\code{eye(n)} make matrices, and \code{zeros(n)} and \code{ones(n)} vectors;
\code{linspace(a, b, n)} gives n equally spaced points and \code{rand(m, n)}
uniform random numbers. Arrays can be joined as blocks inside a literal.

\begin{axon}
x = linspace(0.0, 1.0, 5)
println(x')
I = eye(2)
M = [I, 2 * I; zeros(1, 4)]      # blocks: a 3 by 4 matrix
[m, n] = size(M)
println(m, n, length(M))
w = [x; 2.0]                     # append to a vector
println(length(w), w[end])
\end{axon}
\begin{axonout}
1(*@\texttimes@*)5 matrix{real}
  0.0  0.25  0.5  0.75  1.0
3 4 12
6 2.0
\end{axonout}

\paragraph{Functions of arrays.} Elementary functions such as \code{sin},
\code{exp} and \code{sqrt} apply to each element. \code{sum}, \code{prod},
\code{mean}, \code{min} and \code{max} reduce a whole array to a number;
\code{sum(A, 1)} sums each column and \code{sum(A, 2)} each row.

\begin{axon}
t = linspace(0.0, 1.0, 5)
println(sum(t), mean(t), max(exp(t)))
B = [1, 2, 3; 4, 5, 6]
println(sum(B, 1))
println(sum(B, 2))
println(sort([3; 1; 2])')
println(sortperm([30; 10; 20])')
\end{axon}
\begin{axonout}
2.5 0.5 2.71828182845905
1(*@\texttimes@*)3 matrix{int}
  5  7  9
2(*@\texttimes@*)1 matrix{int}
   6
  15
1(*@\texttimes@*)3 matrix{int}
  1  2  3
1(*@\texttimes@*)3 matrix{int}
  2  3  1
\end{axonout}

\paragraph{Broadcasting.} When the two operands of an elementwise operation
differ in shape, a dimension of size 1 is stretched to match the other: a
scalar combines with every element, and a row combines with every row. So
subtracting the column means from a matrix is one line:

\begin{axon}
X = [1.0, 10.0; 3.0, 30.0]
println(X - mean(X, 1))
\end{axon}
\begin{axonout}
2(*@\texttimes@*)2 matrix{real}
  -1.0  -10.0
   1.0   10.0
\end{axonout}

\paragraph{Solving.} \code{A \ b} solves $Ax=b$, by LU factorisation for a
square matrix and by least squares otherwise; it reports a singular matrix
as an error, \code{SingularMatrix}. \code{A ^ k} is a matrix power. Chapter~\ref{ch:num}
has the rest of linear algebra.

\section{Indexing and ranges}

Indices start at 1. Inside the brackets, \code{end} is the last index, \code{:}
selects a whole row or column, and a range or a vector of indices selects
several elements. A \code{bool} array of the same size selects where it is
\code{true}.

\begin{axon}
w = [10, 20, 30, 40, 50]
println(w[2], w[end], w[end-1])
println(w[2:4])
A = [1, 2, 3; 4, 5, 6; 7, 8, 9]
println(A[2, :])        # a row, as a vector
println(A[:, end]')     # the last column
println(A[A > 5]')      # a mask: the elements above 5
\end{axon}
\begin{axonout}
20 50 40
3-vector{int}
  20
  30
  40
3-vector{int}
  4
  5
  6
1(*@\texttimes@*)3 matrix{int}
  3  6  9
1(*@\texttimes@*)4 matrix{int}
  7  8  6  9
\end{axonout}

A mask selects in column-major order, the order in which a matrix is stored:
down the first column, then the second, and so on.

\paragraph{Ranges.} \code{a:b} counts from a to b in steps of 1 and
\code{a:h:b} in steps of h, which may be negative or fractional. A range is lazy:
it stores only its ends and step, and becomes a vector where one is needed.

\begin{axon}
for k in 1:2:7
  print(k, "")
end
println()
println(length(0:0.1:1), sum(5:-1:1))
\end{axon}
\begin{axonout}
1 3 5 7
11 15
\end{axonout}

Note the order: start, step, stop, as in MATLAB and Julia. The built-in
\code{range} takes the same numbers in reading order, start, stop, step, and
makes the same range; the step is 1 when left out, and \code{range(n)} is
\code{1:n}, the indices of something of length n. Put a range in brackets to
see its elements.

\begin{axon}
r = range(0, 1, 0.25)
println(r)
println([r]')
println(isequal([range(1, 9, 2)], [1:2:9]))
\end{axon}
\begin{axonout}
0.0:0.25:1.0
1(*@\texttimes@*)5 matrix{real}
  0.0  0.25  0.5  0.75  1.0
true
\end{axonout}

\paragraph{Assignment.} Assigning to an element or a slice changes only that
variable. A scalar on the right is copied into every selected element.

\begin{axon}
A = zeros(3, 3)
A[2, :] = [1.0, 2.0, 3.0]
A[:, 1] = 9.0
println(A)
\end{axon}
\begin{axonout}
3(*@\texttimes@*)3 matrix{real}
  9.0  0.0  0.0
  9.0  2.0  3.0
  9.0  0.0  0.0
\end{axonout}

An index outside the array is an \code{IndexError}: reading never wraps around
and assigning never grows an array. To add elements, use \code{append(v, x)} or
a literal such as \code{[v; x]}, which return a new array.

\section{Strings}

Strings are written in double quotes, joined with \code{+}, and indexed by
character. \code{string(x)} gives any value's display form as a string.

\begin{axon}
name = "Axon"
println("Hello, " + name + "! " + string(length(name)) + " letters")
println(name[1], name[2:3], "tab\there", "quote: \"q\"")
for c in "abc"
  print(c, "")
end
println()
\end{axon}
\begin{axonout}
Hello, Axon! 4 letters
A xo tab	here quote: "q"
a b c
\end{axonout}

Strings compare with \code{==} and order with \code{<}. They cannot be changed
in place: build a new one instead. Chapter~\ref{ch:text} covers strings in
full: searching, changing, splitting and reading numbers from text.

\section{Printing and display}\label{sec:printing}

\code{println(a, b, ...)} writes its arguments separated by spaces and ends the
line; \code{print} does the same without ending it, so a line can be built up in
a loop. \code{println()} on its own ends the line.

\begin{axon}
println("x =", 3, 2.5, true)
for k in 1:3
  print(string(k) + " ")
end
println()
println("Total: " + string(42))
\end{axon}
\begin{axonout}
x = 3 2.5 true
1 2 3
Total: 42
\end{axonout}

\paragraph{Choosing the digits.} \code{format(x, d)} gives \code{x} as a string
with exactly \code{d} digits after the point, for building text. The
\code{fmt} module changes how values are displayed from then on, and formats
numbers and arrays with more options (Section~\ref{sec:fmt}):

\begin{itemize}
\item \code{fmt.set_decimals(n)} makes every displayed real show \code{n}
  decimals; \code{fmt.set_decimals()} goes back to significant digits.
\item \code{fmt.set_size(rows, cols)} sets how large an array is shown in full.
  By default, arrays above 10 rows or 8 columns show their first and last rows
  and columns only; \code{fmt.set_size()} restores that.
\item \code{fmt.set_header(false)} leaves out the header line of displayed
  arrays, such as \code|2x3 matrix{int}|, so only the rows are
  printed; \code{fmt.set_header()} brings it back.
\item \code{fmt.set_sci(small, large)} sets when reals switch to scientific
  notation: below \code{small} or from \code{large} in size, instead of below
  \code{1e-4} or from \code{1e15}. \code{fmt.set_sci(1e-6, 1e9)} keeps
  \code{0.00001} as it is but writes \code{2e9} as \code{2.0e9};
  \code{fmt.set_sci()} restores the defaults.
\end{itemize}

\begin{axon}
import fmt
println("pi is about " + format(pi, 2))
fmt.set_decimals(3)
println(1 / 3, [1.0, 22.5])
fmt.set_decimals()
fmt.set_header(false)
println([1, 2; 3, 4])
fmt.set_header()
\end{axon}
\begin{axonout}
pi is about 3.14
0.333 1(*@\texttimes@*)2 matrix{real}
  1.000  22.500
1  2
3  4
\end{axonout}

\section{Control flow}

\code{if}, \code{for} and \code{while} close with \code{end}. Conditions must be
\code{true} or \code{false}: there is no \emph{truthiness} for numbers, strings or
empty collections, and a \code{bool} array is an error (use \code{all} or
\code{any}). \code{and}, \code{or} and \code{not} are words.

\begin{axon}
total = 0
n = 1
while total < 20
  total = total + n
  n = n + 1
end
if total > 30
  println("big")
elseif total > 15 and n < 10
  println("medium", total)
else
  println("small")
end
for x in [3, 1, 4, 1, 5]
  if x == 1
    continue        # skip to the next element
  end
  if x > 4
    break           # leave the loop
  end
  print(x, "")
end
println()
\end{axon}
\begin{axonout}
medium 21
3 4
\end{axonout}

\code{for} walks through a range, the elements of an array (in column-major
order), a list, the characters of a string or the keys of a dictionary. A short
\code{if} fits on one line: \code{if x > 0 return 1 end}.

\section{Functions}

A named function is defined with \code{function} and closes with \code{end}. It
can be called before its definition in the file, because named functions are
bound before the program runs.

\begin{axon}
println(norm2([3.0; 4.0]))

## The Euclidean norm of a vector.
function norm2(v)
  return sqrt(sum(v .^ 2))
end
\end{axon}
\begin{axonout}
5.0
\end{axonout}

\paragraph{Keyword parameters.} Parameters after a \code{;} are keyword
parameters: they have defaults and are passed by name. In calls the \code{;}
is not written, only commas: \code{damped(2.0, rate = 0.1)}.

\begin{axon}
function damped(t; rate = 0.5, amplitude = 1.0)
  return amplitude * exp(-rate * t)
end
println(damped(2.0), damped(2.0, rate = 0.1), damped(2.0, amplitude = 3.0, rate = 0.0))
\end{axon}
\begin{axonout}
0.367879441171442 0.818730753077982 3.0
\end{axonout}

\paragraph{Reading signatures.} This manual and the Language Reference write
a function's parameters the same way:
\code{fmt.matrix(A; decimals: int, digits: int = 6)} means \code{A} is passed
by position and the rest by name. \code{: int} gives a parameter's type, as in
an annotation, and \code{= 6} its default. A keyword shown without a value,
like \code{decimals} here, is \emph{not set} unless you give it (its value is
\code{nothing}), and the function then does what its description says: here,
numbers show their significant digits. Not set is different from zero:
\code{decimals = 0} asks for no digits after the point.

\paragraph{Optional positional parameters.} A positional parameter with a
default may be left out. Every positional parameter after it needs a default
too. A default of \code{nothing}, tested in the body, gives one function two
ways to be called, as \code{plot.bar(y)} and \code{plot.bar(x, y)} do.

\begin{axon}
function area(w, h = w)          # a square unless h is given
  return w * h
end
function span(a, b = nothing)
  if typeof(b) == "nothing"
    return a
  end
  return b - a
end
println(area(3), area(3, 2), span(4), span(4, 9))
\end{axon}
\begin{axonout}
9 6 4 5
\end{axonout}

\paragraph{Several results.} \code{return a, b} returns two values, which the
caller unpacks with brackets on the left of \code{=}.

\begin{axon}
function minmax(v)
  return min(v), max(v)
end
[lo, hi] = minmax([3, 1, 4, 1, 5])
println(lo, hi)
\end{axon}
\begin{axonout}
1 5
\end{axonout}

\paragraph{Functions as values.} \code{function(x) ... end} without a name is an
anonymous function (a \emph{lambda}), a value like any other. It captures the
variables around it by reference, so it sees later changes to them and can
change them.

\begin{axon}
k = 3.0
scale = function(x) return k * x end
println(scale(2.0))
k = 10.0
println(scale(2.0))

count = 0
tick = function() count = count + 1 end
tick(); tick()
println(count)

apply_twice = function(f, x) return f(f(x)) end
println(apply_twice(scale, 1.0))
\end{axon}
\begin{axonout}
6.0
20.0
2
100.0
\end{axonout}

\paragraph{Scope.} Variables first assigned in a function are local to it. A
named top-level function can read the program's variables but needs
\code{global x} to assign one. \code{local x} makes a fresh variable in the
current block. A function may call itself, up to 10\,000 calls deep.

\begin{axon}
counter = 0
function bump()
  global counter
  counter = counter + 1
end
bump(); bump()
println(counter)

function fact(n)
  if n <= 1 return 1 end
  return n * fact(n - 1)
end
println(fact(20), factorial(20))
\end{axon}
\begin{axonout}
2
2432902008176640000 2432902008176640000
\end{axonout}

\section{Lists, records and dictionaries}

Three kinds of collection hold values of any type.

\paragraph{Lists} are written in braces and index like vectors.
\code{append(l, x)} returns a new list with \code{x} added.

\paragraph{Records} have named fields, read and set with a dot. Assigning a new
field adds it.

\paragraph{Dictionaries} map keys (strings or numbers) to values. \code{d[key]}
reads and sets an entry; a missing key is a \code{KeyError}, which
\code{get(d, key, default)} avoids. Entries keep the order they were added in.

\begin{axon}
items = {1, "two", 3.0}
items = append(items, [4; 5])
println(length(items), items[2])

p = {a: 1.0, b: 0.1}
p.c = 0.075
println(p, p.b)

d = dict("Km" => 0.5, "Vm" => 10.0)
d["kcat"] = 2.0
println(d["Km"], keys(d), haskey(d, "Vm"), get(d, "Ki", 0.0))
for k in d
  print(k, "")
end
println()
\end{axon}
\begin{axonout}
4 two
{a: 1.0, b: 0.1, c: 0.075} 0.1
0.5 {"Km", "Vm", "kcat"} true 0.0
Km Vm kcat
\end{axonout}

\paragraph{Declared records.} A \code{record} definition fixes a record's fields
and their types. A value is made by calling the type with every field, by
position or by name, and each field is checked.

\begin{axon}
record Species
  name: string
  init: real
end
s = Species("glucose", 10)        # 10 is widened to 10.0
t = Species(name = "ATP", init = 2.5)
println(s, typeof(t), t.init)
\end{axon}
\begin{axonout}
{name: "glucose", init: 10.0} Species 2.5
\end{axonout}

\section{Errors}\label{sec:errors}

An error has a \code{kind}, such as \code{"IndexError"}, a \code{message} and a
\code{trace}. \code{try ... catch err ... end} catches errors raised in the
\code{try} block, including those raised in functions it calls.
\code{error(kind, message)} raises one, \code{error(message)} raises a
\code{UserError}, and \code{error(err)} passes on an error you caught but do not
handle.

\begin{axon}
function checked_sqrt(x)
  if x < 0 error("DomainError", "negative input: " + string(x)) end
  return sqrt(x)
end
try
  checked_sqrt(-1.0)
catch err
  println(err.kind, "-", err.message)
end

# Handle one kind, pass on the rest.
try
  x = [1.0, 2.0; 2.0, 4.0] \ [1.0; 1.0]
catch err
  if err.kind != "SingularMatrix"
    error(err)
  end
  println("singular, as expected")
end
\end{axon}
\begin{axonout}
DomainError - negative input: -1.0
singular, as expected
\end{axonout}

The kinds an Axon program can meet are listed in Table~\ref{tab:errors}.

\begin{longtable}{@{}lp{0.68\linewidth}@{}}
\caption{Kinds of error.}\label{tab:errors}\\
\toprule
Kind & Raised when \\
\midrule
\endfirsthead
\toprule
Kind & Raised when \\
\midrule
\endhead
\texttt{ArgumentError} & Wrong number of arguments, unknown keyword, invalid
  argument value \\
\texttt{DimensionError} & Incompatible shapes in an operation, literal or
  assignment \\
\texttt{DomainError} & An argument outside a function's domain with no NaN
  result, such as \code{factorial(-1)} \\
\texttt{FieldError} & A missing field of a record \\
\texttt{ForeignError} & A native library or symbol cannot be found, or a native
  call fails \\
\texttt{ImportError} & A module cannot be found, or imports form a cycle \\
\texttt{IndexError} & An index is out of bounds or not an integer \\
\texttt{IOError} & A file cannot be read, written or deleted \\
\texttt{KeyError} & A dictionary key is missing \\
\texttt{NameError} & A name is read before it is bound \\
\texttt{OverflowError} & Integer arithmetic exceeds 64 bits \\
\texttt{RecursionError} & Calls nest more than 10\,000 deep \\
\texttt{SingularMatrix} & A solve or inverse meets a singular matrix \\
\texttt{SolverError} & A numerical solver fails, such as an ODE whose step
  becomes too small \\
\texttt{SyntaxError} & The program text cannot be read; nothing runs \\
\texttt{TypeError} & A value of the wrong type for an operation, annotation or
  condition \\
\texttt{UserError} & Raised by \code{error("text")} \\
\bottomrule
\end{longtable}

\section{Modules}\label{sec:modules}

\code{import m} loads a module and binds it as \code{m}, so its functions are
\code{m.name}; \code{import m as x} binds it under another name, and
\code{from m import f, g} binds chosen functions directly. A module is loaded
once, however many times it is imported.

\begin{axon}
import linalg
from linalg import eig, det
import random as rnd
println(linalg.det([2.0, 1.0; 1.0, 2.0]), det(eye(3)))
println(eig([2.0, 0.0; 0.0, 3.0])')
rnd.seed(1)
\end{axon}
\begin{axonout}
3.0 1.0
1(*@\texttimes@*)2 matrix{real}
  2.0  3.0
\end{axonout}

\paragraph{Writing a module.} Any \code{.sci} file is a module named after the
file; a folder of files is a namespace, so \texttt{mymodels/kinetics.sci} is
\code{mymodels.kinetics}. Only definitions marked \code{export} are visible to
programs that import it. A \code{##} block at the top of the file is the
module's documentation.

\ex{mymodels/kinetics.sci}
\ex{use_module.sci}
\begin{axonout}
0.5 8.0
0.5
\end{axonout}

\paragraph{Where modules are found.} \code{import m} looks in the importing
file's own folder first, then in the packages listed in the nearest
\texttt{package.toml} file, then in the standard library. The standard library
modules are listed in Table~\ref{tab:modules}.

\begin{table}[htbp]
\centering
\caption{The standard modules.}\label{tab:modules}
\begin{tabular}{@{}lp{0.7\linewidth}@{}}
\toprule
Module & Contents \\
\midrule
\texttt{linalg} & Solving, determinants, inverses, eigenvalues, SVD, QR, LU,
  Cholesky, rank, null space, norms, matrix exponential \\
\texttt{ode} & \code{ode.solve}: ordinary differential equations, explicit
  (RK45) and stiff (BDF), with events and dense output \\
\texttt{plot} & Line, scatter, bar, histogram, surface and 3D plots, subplots;
  PNG, SVG and PDF \\
\texttt{random} & Seeded uniform, normal and integer random numbers \\
\texttt{time} & Clocks and timing code \\
\texttt{io} & Reading and writing files: text, lines, CSV and JSON (Chapter~\ref{ch:files}) \\
\texttt{fmt} & Display settings, and numbers and arrays as text
  (Section~\ref{sec:fmt}) \\
\texttt{graphics} & The drawing layer under \texttt{plot} \\
\bottomrule
\end{tabular}
\end{table}

% ===================================================================
\chapter{Numerical computing}\label{ch:num}

\section{Linear algebra}

The core operators handle the everyday cases: \code{A * B} multiplies through
an optimised BLAS library, \code{A \ b} solves a square system by LU
factorisation and an over- or under-determined one by least squares,
\code{A / B} solves from the right, and \code{A ^ k} is a matrix power. A
singular system raises \code{SingularMatrix}. \code{A'} is the transpose
(without conjugating complex entries), and so is \code{tr(A)};
\code{adjoint(A)} is the conjugate transpose and \code{dot(u, v)} the inner
product with \code{u} conjugated, for complex data.

The \code{linalg} module adds the decompositions. The table also lists
\code{tr} and \code{trace}, which are built in and need no \code{import}:

\begin{longtable}{@{}p{0.3\linewidth}p{0.64\linewidth}@{}}
\toprule
Function & Result \\
\midrule
\endhead
\code{tr(A)} & The transpose, the same as \code{A'} (built in) \\
\code{trace(A)} & The sum of the diagonal of a square matrix (built in) \\
\code{solve(A, b)} & \code{A \ b} for square \code{A} \\
\code{lstsq(A, b)} & The minimum-norm least-squares solution, for any shape \\
\code{det(A)} & The determinant \\
\code{inv(A)} & The inverse \\
\code{eig(A; vectors: bool = false)} & The eigenvalues; \code{eig(A, vectors = true)} also returns the
  eigenvectors \\
\code{svd(A)} & \code{U, s, V} with \code{A = U*diag(s)*V'}, \code{s} in
  decreasing order \\
\code{qr(A)} & \code{Q, R}, the thin QR factorisation \\
\code{lu(A)} & \code{L, U} with \code{A = L*U} (\code{L} row-permuted); works
  on sparse matrices too \\
\code{chol(A)} & Upper-triangular \code{R} with \code{A = R'*R} \\
\code{rank(A; tol: real)} & The numerical rank \\
\code{null(A; tol: real)} & An orthonormal basis of the null space \\
\code{norm(x; p = 2)} & Vector p-norms; for a matrix the 2-norm, \code{p = 1},
  \code{"inf"} or \code{"fro"} \\
\code{expm(A)} & The matrix exponential \\
\bottomrule
\end{longtable}

\ex{linear_algebra.sci}
\begin{axonout}
x = 1(*@\texttimes@*)3 matrix{real}
  -0.302326  0.767442  0.72093
residual: 0.0
det: 43.0
eigenvalues: 1(*@\texttimes@*)3 matrix{real}
  1.8549  3.47602  6.66908
R'R matches A: true
singular values: 1(*@\texttimes@*)3 matrix{real}
  6.66908  3.47602  1.8549
least squares: 1(*@\texttimes@*)2 matrix{real}
  3.5  1.4
\end{axonout}

Array entries are shown with 6 significant digits; \code{format} or
\code{fmt.matrix} give more.

\section{Sparse matrices}

A sparse matrix stores only its nonzero entries, which makes systems with
thousands of unknowns cheap when each equation involves only a few of them.
\code{sparse(i, j, v, m, n)} builds one from index vectors and values (summing
repeated entries); \code{speye(n)} is the sparse identity, \code{nnz(S)} counts
the stored entries and \code{full(S)} converts to a dense matrix. Sums, products
and products with scalars stay sparse, and \code{S \ b} solves with a sparse LU
factorisation built into the interpreter.

\ex{sparse_poisson.sci}
\begin{axonout}
stored entries: 5998 of 4000000
u(0.5) about 0.125000  exact 0.125000
\end{axonout}

\section{Differential equations}\label{sec:ode}

\code{ode.solve(f, y0, tspan)} solves the system $y' = f(t, y)$ from
\code{tspan[1]} to \code{tspan[2]}, starting from \code{y0}. The function
\code{f(t, y)} receives a real time and a real vector and returns the
derivative as a vector of the same length (or a number, for one equation).

\ex{lotka.sci}
\begin{axonout}
status: success  steps: 343
prey at t = 30: 40.164
state at t = 12.5: 2-vector{real}
  26.9341
  3.44952
\end{axonout}

\fig{lotka.pdf}{0.8}{The Lotka--Volterra populations, from \texttt{lotka.sci}.}

\paragraph{The result.} \code{ode.solve} returns a record:

\begin{longtable}{@{}lp{0.68\linewidth}@{}}
\toprule
Field & Contents \\
\midrule
\endhead
\code{t} & The times reported: every step, or the times asked for with
  \code{t_eval} \\
\code{y} & A matrix with one column per reported time, so \code{sol.y[1, :]}
  is the first component over time \\
\code{at} & A function giving the state at any time in the solved range:
  \code{sol.at(2.5)}, or a matrix for a vector of times \\
\code{status} & \code{"success"}, or \code{"event"} when an event stopped the
  integration \\
\code{t_events}, \code{y_events} & For each event function, the times and states
  at which it occurred \\
\code{names} & The component names \\
\code{nsteps}, \code{nfev} & Accepted steps and evaluations of \code{f} \\
\code{method} & The method used \\
\bottomrule
\end{longtable}

\paragraph{Options.} Keyword arguments control the solution:

\begin{longtable}{@{}p{0.3\linewidth}p{0.64\linewidth}@{}}
\toprule
Keyword & Meaning (default) \\
\midrule
\endhead
\code{method} & \code{"rk45"}, explicit Dormand--Prince (default), or
  \code{"bdf"} for stiff problems \\
\code{rtol}, \code{atol} & Relative and absolute tolerances (\code{1e-6},
  \code{1e-9}); \code{atol} may be a vector, one tolerance per
  component, when the components differ greatly in size \\
\code{t_eval} & The times at which to report the solution \\
\code{events} & An event function \code{g(t, y)}, or a list of them \\
\code{terminal} & Whether an event stops the integration (\code{false}) \\
\code{direction} & Which crossings count: \code{0} both, \code{1} rising,
  \code{-1} falling \\
\code{names} & Names for the components, a list of strings \\
\code{max_step}, \code{first_step}, \code{max_steps} & Step-size limits, and
  the most steps before giving up (100\,000) \\
\bottomrule
\end{longtable}

\paragraph{Stiff problems and events.} A problem is \emph{stiff} when it mixes
very fast and very slow processes; an explicit method then needs tiny steps
throughout. \code{method = "bdf"} uses the variable-order BDF method of SUNDIALS
CVODE, which takes large steps on such problems. An \emph{event} is a function
of \code{t} and \code{y} whose zero crossing matters, such as a ball reaching
the ground; with \code{terminal = true} the integration stops there.

\ex{stiff_events.sci}
\begin{axonout}
BDF steps: 409
at t = 40: 0.715827 0.0000091855 0.284164
status: event
landed at t = 2.446483  exact: 2.446483
\end{axonout}

\fig{ball.pdf}{0.75}{The ball's height until it lands, from
\texttt{stiff\_events.sci}.}

The stepping loops of both methods run as native code, and so does a
right-hand side that the JIT compiles (Chapter~\ref{ch:perf}).

\section{Random numbers}

The prelude's \code{rand} and the \code{random} module draw from one generator.
A program starts from a seed taken from the clock; \code{random.seed(n)} makes
every later draw the same on every run and every platform. Each distribution
takes its parameters, then an optional shape: nothing for a number, \code{n} for
a vector, \code{m, n} for a matrix.

\begin{center}
\begin{tabular}{@{}lp{0.5\linewidth}@{}}
\toprule
Function & Result \\
\midrule
\code{seed(n: int)} & Seeds the generator \\
\code{rand(...)} & Uniform reals in $[0, 1)$ \\
\code{randn(...)} & Standard normal reals \\
\code{uniform(a: real, b: real, ...)} & Uniform reals in $[a, b)$ \\
\code{normal(mu: real, sigma: real, ...)} & Normal reals with mean \code{mu} and
  standard deviation \code{sigma} \\
\code{randint(lo: int, hi: int, ...)} & Integers from \code{lo} to \code{hi} inclusive \\
\bottomrule
\end{tabular}
\end{center}

\ex{random_numbers.sci}
\begin{axonout}
dice: 1(*@\texttimes@*)8 matrix{int}
  1  1  6  6  5  1  5  4
mean: 10.014  sd: 1.986
pi is about 3.1350
\end{axonout}

\section{Timing}

The \code{time} module reads clocks and times code. \code{time.clock()} is a
high-resolution clock for differences; \code{time.elapsed(f, args...)} times
one call; \code{time.bench(f, args...)} repeats the call until the timing is
reliable and returns a record whose \code{best} field is the seconds per call;
\code{time.format(seconds)} writes a time in a fitting unit. An anonymous
function times any expression: \code{time.bench(function() return A * A end)}.
Section~\ref{sec:bench} shows them at work.

% ===================================================================
\chapter{Working with text}\label{ch:text}

Text turns up in every program: labels and titles, names in data files,
messages and reports. This chapter covers how to write strings, take them
apart, search and change them, and put them back together.

\section{Writing strings}

A string is text in double quotes. A backslash starts an escape for a
character that is awkward to type:

\begin{longtable}{@{}lp{0.7\linewidth}@{}}
\toprule
Escape & Character \\
\midrule
\endhead
\code{\n} & A line break \\
\code{\t} & A tab \\
\code{\"} & A double quote \\
\code{\\} & A backslash \\
\code|\u{3B1}| & The character with that hexadecimal code, here the Greek
  letter alpha \\
\bottomrule
\end{longtable}

There are no single-quoted strings, because \code{'} is the transpose
operator. A string in ordinary double quotes cannot run across a line break in
the program; join pieces with \code{+}, which may end a line:

\begin{axon}
note = "First line\n\tthen an indented one, " +
       "and \"quotes\" and \u{3B1} too"
println(note)
\end{axon}
\begin{axonout}
First line
	then an indented one, and "quotes" and (*@$\alpha$@*) too
\end{axonout}

\paragraph{Text across lines.} For longer text, such as a model definition or a
block of data, use a triple-quoted string: it starts with three double quotes,
runs across as many lines as needed, and ends at the next three double quotes.
The text is taken as written, line breaks included, and a \code{"} inside needs
no escape.

\begin{axon}
source = """
    2 A + 3 B -> C; k1*A*B;
    C -> A + B; k2*C;
    k1 = 0.1;
    k2 = 0.05;
    """
print(source)
println(length(split(source, "\n")), "pieces")
\end{axon}
\begin{axonout}
2 A + 3 B -> C; k1*A*B;
C -> A + B; k2*C;
k1 = 0.1;
k2 = 0.05;
5 pieces
\end{axonout}

Two adjustments keep such strings tidy in a program. A line break straight
after the opening quotes is dropped, so the text can start on its own line;
and the indentation the lines share is removed, so the string can be indented
along with the code around it. In the example every line, including the one
with the closing quotes, starts with four spaces, and none of them is in the
string; the string ends with a line break, because the closing quotes are on a
line of their own. Escapes such as \code{\n} and \code|\u{3B1}| still work.

\section{Characters, length and slices}

Indexing picks characters, starting from 1, and a range of indexes gives a
slice; \code{end} is the last position. \code{length} counts characters, not
bytes, so accented letters and other alphabets count one each. Strings cannot
be changed in place: \code{s[1] = "j"} raises \code{TypeError}, so build a new
string instead.

\begin{axon}
s = "Bacillus"
println(s[1], s[end], s[2:4], s[end - 2:end], length(s))
println(length("na\u{EF}ve"))
for c in "abc"
  print(c, "")
end
println()
t = "M" + s[2:end]       # a new string; s is unchanged
println(t, s)
\end{axon}
\begin{axonout}
B s aci lus 8
5
a b c
Macillus Bacillus
\end{axonout}

\section{Building strings}

\begin{itemize}
\item \code{+} joins two strings. Numbers must be converted first:
  \code{"n = " + string(n)}; \code{"n = " + n} raises \code{TypeError}.
\item \code{string(x)} gives any value as text, as \code{println} would show
  it, and \code{format(x, d)} a number with exactly \code{d} decimals.
\item \code{repeat(s, n)} repeats a string, which is handy for rules and
  padding, and \code{join(parts, sep)} joins a list of strings with a separator.
\item \code{println(a, b, c)} prints its values separated by single spaces;
  to control the spacing exactly, build one string with \code{+} and print
  that.
\end{itemize}

\begin{axon}
x = 2.0 / 3.0
println("x = " + string(x) + ", or about " + format(x, 3))
println(repeat("=", 20))
println(join({"Na", "K", "Cl"}, ", "), "|", join({"a", "b"}))
\end{axon}
\begin{axonout}
x = 0.666666666666667, or about 0.667
====================
Na, K, Cl | ab
\end{axonout}

\section{Comparing and sorting}

\code{==} and \code{!=} compare whole strings, and \code{<}, \code{<=},
\code{>} and \code{>=} compare them by character code, which is alphabetical
order except that all capitals come before all small letters. \code{sort} on a
list of strings uses the same order. To ignore case, compare or sort
lowered copies:

\begin{axon}
names = {"pear", "Apple", "fig"}
println(sort(names), "Zebra" < "apple")
println(lower("Apple") < lower("apple") or lower("Apple") == lower("apple"))
\end{axon}
\begin{axonout}
{"Apple", "fig", "pear"} true
true
\end{axonout}

\section{Searching}

\begin{longtable}{@{}p{0.34\linewidth}p{0.6\linewidth}@{}}
\toprule
Function & Result \\
\midrule
\endhead
\code{contains(s: string, part: string)} & Whether \code{part} occurs anywhere in \code{s} \\
\code{startswith(s: string, part: string)} & Whether \code{s} begins with \code{part} \\
\code{endswith(s: string, part: string)} & Whether \code{s} ends with \code{part} \\
\code{find(s: string, part: string)} & The position of the first \code{part} in \code{s}, or 0
  when there is none \\
\code{find(s, part, start: int)} & The same, looking from position \code{start} \\
\bottomrule
\end{longtable}

Positions from \code{find} are character positions, so they work directly as
indexes. Searching again from just after the last match finds every
occurrence:

\begin{axon}
s = "abracadabra"
println(contains(s, "cad"), startswith(s, "abr"), endswith(s, ".csv"))
k = find(s, "a")
while k > 0
  print(k, "")
  k = find(s, "a", k + 1)
end
println()
\end{axon}
\begin{axonout}
true true false
1 4 6 8 11
\end{axonout}

\section{Changing strings}

These functions return a new string and leave the original alone:

\begin{longtable}{@{}p{0.34\linewidth}p{0.6\linewidth}@{}}
\toprule
Function & Result \\
\midrule
\endhead
\code{upper(s: string)}, \code{lower(s: string)} & \code{s} in capitals or small letters, for
  any alphabet \\
\code{trim(s: string)} & \code{s} without spaces, tabs and line breaks at either end \\
\code{ltrim(s: string)}, \code{rtrim(s: string)} & The same at the start only, or at the end
  only \\
\code{replace(s: string, old: string, new: string)} & \code{s} with every \code{old} replaced by
  \code{new} \\
\bottomrule
\end{longtable}

\begin{axon}
s = "   Sample 7: OD = 0.42   "
println("[" + trim(s) + "]", "[" + ltrim(s) + "]", "[" + rtrim(s) + "]")
println(upper("ph"), lower("NaCl"), replace("2024-03-15", "-", "/"))
\end{axon}
\begin{axonout}
[Sample 7: OD = 0.42] [Sample 7: OD = 0.42   ] [   Sample 7: OD = 0.42]
PH nacl 2024/03/15
\end{axonout}

\section{Splitting text and reading numbers}

\code{split(s)} breaks a string into words at runs of spaces, tabs and line
breaks; \code{split(s, sep)} breaks it at every \code{sep}, keeping empty
pieces, as a list of strings. \code{int(s)} and \code{real(s)} read the number
written in a string, ignoring spaces around it, and raise
\code{ArgumentError} for anything else. Chapter~\ref{ch:files} uses them to
read data files line by line.

\begin{axon}
words = split("  the  quick brown\tfox ")
println(words, length(words))
fields = split("12, 4.5, , 7", ",")
println(fields)
println(int(trim(fields[1])) + 1, real(fields[2]) * 2)
\end{axon}
\begin{axonout}
{"the", "quick", "brown", "fox"} 4
{"12", " 4.5", " ", " 7"}
13 9.0
\end{axonout}

\section{Functions from text}\label{sec:parsefn}

A formula can arrive as text: typed by the user, or read from a settings file.
\code{parse_function} turns it into a function that can be called, passed to
\code{ode.solve} or plotted like any other:

\begin{axon}
f = parse_function("x, y", "2*sin(x) + y^2")
println(f(1.0, 2.0))
g = parse_function("function(x)\n  if x < 0\n    return -x\n  end\n  return x^2\nend")
println(g(-3), g(3))
\end{axon}
\begin{axonout}
5.68294196961579
3 9
\end{axonout}

The first form takes the parameter names and the expression the function
returns; the second takes a whole \code{function(...) ... end}, for formulas
that need more than one line. A few things to know:

\begin{itemize}
\item \textbf{What a formula can use.} Its parameters, and the built-in
  functions and constants (\code{sin}, \code{exp}, \code{pi}, \ldots). It does
  not see the program's own variables or any module, so a formula from a file
  cannot reach, say, \code{io.delete_file}. Pass any other value in as a
  parameter.
\item \textbf{Mistakes are caught at once.} An unknown name raises
  \code{NameError} and an expression that does not parse \code{SyntaxError},
  both saying where in the formula the problem is, when
  \code{parse_function} is called rather than later.
\item \textbf{Whole vectors.} To evaluate a formula over a vector of points in
  one call, as for a plot, write it with the elementwise operators: \code{.^},
  \code{.*} and \code{./}.
\item \textbf{Speed.} The result is compiled by the JIT at its first call,
  like any other function.
\end{itemize}

\ex{formula_plot.sci}
\begin{axonout}
NameError: parse_function: unknown name "sinn" at character 5 of "2 * sinn(x)"
SyntaxError: parse_function: the expression "(x + 1" is incomplete (an unclosed bracket, or an operator with nothing after it?)
\end{axonout}

\fig{formulas.pdf}{0.75}{Formulas given as text, from \texttt{formula\_plot.sci}.}

\section{Formatting numbers and arrays}\label{sec:fmt}

\code{println} shows numbers and arrays in a standard form. For reports,
tables and files you often want something else: a fixed number of decimals,
no header line, commas between values. The \code{fmt} module's
\code{number} and \code{matrix} return the text, with keyword options, and
leave the display settings alone:

\begin{longtable}{@{}p{0.3\linewidth}p{0.64\linewidth}@{}}
\toprule
Keyword & Effect \\
\midrule
\endhead
\code{decimals: int} & Exactly that many digits after the point \\
\code{digits: int} & Up to that many significant digits (\code{matrix} uses 6,
  as the display does) \\
\code{sci: bool} & \code{true} for scientific notation, with \code{decimals}
  digits after the point \\
\code{width: int} & \code{number} only: pad on the left to that many
  characters \\
\code{header: bool} & \code{matrix} only: \code{true} adds the
  \code|2x3 matrix{real}| line \\
\code{sep: string} & \code{matrix} only: what goes between values (two
  spaces) \\
\code{align: bool} & \code{matrix} only: \code{false} drops the padding that
  lines up the columns \\
\bottomrule
\end{longtable}

\begin{axon}
import fmt
A = [1.5, -22.25; 3.0, 4.125]
println(fmt.matrix(A))
println(fmt.matrix(A, decimals = 2))
println(fmt.matrix(A, decimals = 1, sep = ", ", align = false))
println(fmt.number(pi, digits = 4), fmt.number(12345.678, sci = true, decimals = 2))
println("[" + fmt.number(42, width = 6) + "]")
\end{axon}
\begin{axonout}
1.5  -22.25
3.0   4.125
1.50  -22.25
3.00    4.12
1.5, -22.2
3.0, 4.1
3.142 1.23e4
[    42]
\end{axonout}

A few points. \code{fmt.matrix} always includes every row and column; the
display limits of \code{fmt.set_size} apply only to printed arrays. A vector
comes out as a column, one value per line; transpose it (\code{v'}) for a
row. Rounding to a number of decimals rounds the exact binary value, ties to
even, so \code{4.125} with two decimals is \code{4.12}; and
\code{format(2.675, 2)} is \code{2.67}, because the double nearest 2.675 is
just below it.

\section{An example: tidying names and a report}

Hand-typed data is rarely consistent. This program brings a list of messy
species names to one spelling each, counts them, and prints an aligned table,
using most of the functions in this chapter.

\ex{text_report.sci}
\begin{axonout}
Organism           Count
------------------------
Bacillus subtilis  2
E. coli            2
Pseudomonas        1
Yeast              1
Bacillus subtilis matches; "i" first at 4
E. coli matches; "i" first at 7
E. coli MG1655
\end{axonout}

The \code{pad} function shows a common pattern: a column lines up when each
entry is padded with \code{repeat(" ", width - length(s))} to the width of the
longest one, and each line is built as one string so that \code{println}
adds no spaces of its own.

% ===================================================================
\chapter{Files and data}\label{ch:files}

Data comes in files and results go out to them. The \code{io} module reads
and writes whole text files, lists of lines, tables of numbers in CSV form, and
any Axon value as JSON. A handful of built-in string functions turn text into
numbers. File names are relative to the folder the program runs in (in the
IDE, the folder of the program's file); a file that cannot be read or written
raises \code{IOError}.

\begin{longtable}{@{}p{0.36\linewidth}p{0.58\linewidth}@{}}
\toprule
Function & Result \\
\midrule
\endhead
\code{io.read_file(path: string)} & The whole of a text file, as one string \\
\code{io.write_file(path: string, text: string)} & Writes a string, replacing the file \\
\code{io.read_lines(path: string)} & The lines of a text file, as a list of strings \\
\code{io.read_csv(path: string)} & A table of numbers: the record
  \code|{names, data}| \\
\code{io.write_csv(path: string, A, names)} & Writes a matrix as a table, with an
  optional header \\
\code{io.write_json(path: string, value)} & Writes any value (numbers, strings, lists,
  records, vectors, matrices) as JSON \\
\code{io.read_json(path: string)} & The value a JSON file describes \\
\code{io.file_exists(path: string)} & Whether the file exists \\
\code{io.delete_file(path: string)} & Deletes the file; \code{false} if there was none \\
\code{io.read_bytes(path: string)} & The bytes of any file, as an \code{int} vector \\
\bottomrule
\end{longtable}

\section{Tables of numbers: CSV}

A CSV file holds a table as text: one row per line, the values separated by
commas, and often a first line of column names. Spreadsheets, instruments and
most other programs can read and write it, which makes it the usual way to
move measurements in and out.

\begin{itemize}
\item \code{io.write_csv(path, A, names)} writes the matrix \code{A}, one row
  per line, under a header line of column names; leave \code{names} out for no
  header. A vector is written as a single column. Numbers are written so that
  they read back to exactly the same value.
\item \code{io.read_csv(path)} returns a record with two fields:
  \code{data}, a \code{real} matrix with one row per line, and \code{names},
  the column names, or \code{nothing} when the file has no header. Axon treats
  the first line as a header when any field in it is not a number.
\end{itemize}

\ex{measurements.sci}
\begin{axonout}
minutes,celsius
0,81.1
2,71.6
4,61.8
{"minutes", "celsius"}
16 rows, 2 columns
first: 81.1 last: 25.2
\end{axonout}

\fig{cooling.pdf}{0.6}{Readings written to a CSV file and read back, from
\texttt{measurements.sci}.}

A few points about the example. \code{[t, temp]} puts the two column vectors
side by side as a 16\texttimes 2 matrix, the shape \code{write_csv} expects.
After reading, \code{table.data[:, 1]} is the first column, and the header's
names label the axes.

\paragraph{Details.}
\begin{itemize}
\item \textbf{Other separators.} Files from some countries and programs use
  semicolons or tabs: \code{io.read_csv(path, ";")} and
  \code{io.read_csv(path, "\t")} read them, and
  \code{io.write_csv(path, A, names, ";")} writes them (pass \code{nothing}
  for \code{names} to write no header).
\item \textbf{Missing values.} An empty field is read as NaN, and NaN is
  written as \code{NaN}. Test for it with \code{x != x}, which is true only for
  NaN.
\item \textbf{Quoted fields.} A field in double quotes may contain the
  separator, as in \code{"conc, mM"}; \code{write_csv} adds the quotes when a
  name needs them.
\item \textbf{Problems} are reported with the line: a row with the wrong number
  of fields, or a field that is not a number, raises \code{IOError} naming the
  line and field. Blank lines are skipped.
\item \textbf{Text columns.} \code{read_csv} reads numbers only. For a file
  with columns of text, read it with \code{io.read_lines} and take each line
  apart with \code{split}, as the next section shows.
\end{itemize}

\section{Text files, line by line}

\code{io.read_lines(path)} gives the lines of a text file as a list of
strings, without their line endings, whether the file was written on Windows
or not. These built-in functions then take each line apart
(Chapter~\ref{ch:text} has the rest of the string functions):

\begin{longtable}{@{}p{0.3\linewidth}p{0.64\linewidth}@{}}
\toprule
Function & Result \\
\midrule
\endhead
\code{split(s: string)} & The words of \code{s}: the pieces between runs of spaces,
  tabs and line breaks \\
\code{split(s, sep: string)} & The pieces between each \code{sep}, empty ones
  included: \code{split("a,,b", ",")} is \code|{"a", "", "b"}| \\
\code{trim(s: string)} & \code{s} without spaces, tabs and line breaks at either end \\
\code{join(parts: list, sep: string)} & The strings of a list joined, with \code{sep}
  between them \\
\code{real(s: string)}, \code{int(s: string)} & The number written in \code{s}:
  \code{real("2.5e-3")}, \code{int("42")}; anything else raises
  \code{ArgumentError} \\
\bottomrule
\end{longtable}

This program reads a small settings file of \code{key = value} lines, skipping
blank lines and comments, then saves results as JSON:

\ex{settings.sci}
\begin{axonout}
logistic growth - rate: 0.35 capacity: 120
logistic growth vector 4
\end{axonout}

\code{real} reads a number exactly as Axon reads it in a program, correctly
rounded, so \code{real("0.1") == 0.1}. It also accepts \code{NaN},
\code{Inf} and \code{-Inf}, and ignores spaces around the number. Going the
other way, \code{string(x)} gives the text for a value and \code{format(x, d)}
a fixed number of decimals, which \code{join} or \code{+} can assemble into a
line for \code{io.write_file}.

\section{Saving values: JSON}

\code{io.write_json(path, value)} writes almost any Axon value as JSON, a text
format other languages read too, and \code{io.read_json(path)} brings it back:
records come back as records, lists as lists, and arrays of numbers as vectors
and matrices. That makes it the easy way to keep a set of parameters, or a
whole record of results, between runs. Complex numbers, sparse matrices and
functions cannot be written. The Language Reference's ``The io module''
section gives the exact correspondence between JSON and Axon values.

\section{Checking for files}

\begin{axon}
import io
if not io.file_exists("cache.json")
  io.write_json("cache.json", {runs: 0})
end
c = io.read_json("cache.json")
println("runs so far:", c.runs)
println(io.delete_file("cache.json"), io.delete_file("cache.json"))
\end{axon}
\begin{axonout}
runs so far: 0
true false
\end{axonout}

\code{io.delete_file} returns \code{true} when it deleted the file and
\code{false} when there was nothing to delete. Any other failure, such as a
file that cannot be read, raises \code{IOError}, which \code{try} and
\code{catch} can handle like any other error (Section~\ref{sec:errors}).

% ===================================================================
\chapter{Plotting}\label{ch:plot}

The \code{plot} module draws figures. Plotting calls add to the \emph{current
figure}; \code{plot.figure} starts a new one with a given size and title, and
\code{plot.save} writes it to a PNG, SVG or PDF file.

\section{Showing and saving}

\code{plot.show()} displays the current figure: in the IDE's plot pane when the
program runs from the IDE, otherwise in the system's image viewer. You rarely
need to call it, because a figure that has changed since it was last shown or
saved is shown automatically, when the program ends and when
\code{plot.figure()} starts the next figure. A program that draws three figures
therefore shows all three. A program that saves each figure with
\code{plot.save}, and only wants files, opens nothing.

\begin{axon}
import plot
x = linspace(0.0, 2 * pi, 100)
plot.figure(title = "Two waves")
plot.line(x, sin(x), label = "sin")
plot.line(x, cos(x), label = "cos", style = "--")
plot.save("_waves.png")            # .png, .svg or .pdf
plot.save("_waves.pdf")
import io
io.delete_file("_waves.png")
io.delete_file("_waves.pdf")
\end{axon}

Sizes are in points (1/72 inch); the default figure is 640 by 480.
\code{plot.save(path, scale = 3)} sets the pixels per point of a PNG.

\section{Line plots}

\code{plot.line(x, y)} draws a line through the points. Its keywords set the
\code{label} shown in the legend, the \code{color}, the \code{width} in points,
the \code{style} (\code{"-"}, \code{"--"}, \code{":"} or \code{"-."}) and a
\code{marker} at each point (\code{"circle"}, \code{"square"},
\code{"diamond"}, \code{"triangle"}, \code{"plus"} or \code{"cross"}). Series
without a colour take the next of ten colours in turn.

\ex{lines.sci}
\fig{lines.pdf}{0.8}{Line styles and a legend outside the plot, from
\texttt{lines.sci}.}

\code{plot.title}, \code{plot.xlabel} and \code{plot.ylabel} label the figure;
\code{plot.xlim} and \code{plot.ylim} fix the axis ranges (otherwise each is
the data range padded by 5\%); \code{plot.grid(false)} and
\code{plot.legend(false)} hide the grid and the legend. The legend lists the
series that have a label. \code{plot.legend_position} places it:
\code{"best"} (the default, the corner that hides the least data),
\code{"top right"}, \code{"top left"}, \code{"bottom right"},
\code{"bottom left"}, or \code{"outside"}, to the right of the plot.

Colours are names (\code{"red"}, \code{"blue"}, \code{"gray"} and others),
\code{"#rrggbb"} hex codes, or vectors \code{[r, g, b]} with components from 0
to 1.

\section{Scatter plots}

\code{plot.scatter(x, y)} draws a marker at each point, with keywords
\code{label}, \code{color}, \code{size} and \code{marker}.

\ex{scatter_fit.sci}
\begin{axonout}
intercept: 2.198  slope: 1.452
\end{axonout}
\fig{scatter.pdf}{0.72}{Data and a least-squares line, from
\texttt{scatter\_fit.sci}.}

\section{Bar charts}

\code{plot.bar(x, heights)} draws a bar of each height centred on each
position, \code{width} wide (0.8 by default). With the heights alone,
\code{plot.bar(y)}, the bars go at 1, 2, 3, \ldots{} When \code{x} is a list of
strings, the bars go at 1, 2, 3, \ldots{} and the axis is labelled with the
names. To place groups side by side, shift \code{x}, narrow the bars, and name
the groups with \code{plot.xticks(positions, labels)}.

\ex{bars.sci}
\fig{bars.pdf}{1.0}{Three kinds of bar chart in one figure, from
\texttt{bars.sci}.}

\section{Histograms}

\code{plot.hist(x, bins = 10)} counts the values in equal-width bins and draws
one bar per bin; \code{bins} may also be a vector of bin edges.
\code{plot.hist_density} scales the bars so that their total area is 1, for
comparison with a probability density. Both return a record with the
\code{counts}, the bin \code{edges} and the bar \code{heights}.

\ex{histogram.sci}
\begin{axonout}
bins: 30  largest count: 218
sample mean: 4.950
\end{axonout}
\fig{histogram.pdf}{0.75}{A histogram against the normal density, from
\texttt{histogram.sci}.}

\section{Surfaces and 3D plots}

\code{plot.surface(x, y, Z)} draws the surface with height \code{Z[i, j]} above
the point \code{(x[j], y[i])}, coloured by height (\code{colormap = "viridis"} or \code{"gray"}). \code{plot.line3(x, y, z)} draws a curve through
points in space, such as a trajectory, and \code{plot.scatter3(x, y, z)}
markers. They can share a figure, but a figure cannot mix 3D series with 2D
ones. \code{plot.view(azim, elev)} sets the direction you look from, in
degrees, as MATLAB's \texttt{view} does, and \code{plot.zlabel} labels the
vertical axis.

\ex{surface.sci}
\fig{surface.pdf}{0.65}{A surface, from \texttt{surface.sci}.}

\ex{lorenz.sci}
\fig{lorenz.pdf}{0.65}{The Lorenz attractor as a 3D curve, from
\texttt{lorenz.sci}.}

In the IDE, 3D figures appear in an interactive chart that you rotate by
dragging (Section~\ref{sec:plotpane}).

\section{Subplots}

\code{plot.subplot(rows, cols, k)} divides the figure into a grid and makes cell
\code{k} current, so the plotting calls that follow go to it. Cells are
numbered along each row from the top left, as in MATLAB: in a 2 by 2 grid,
cells 1 and 2 are on top. Each cell has its own axes, labels, legend and limits,
2D and 3D cells can be mixed, and the title given to \code{plot.figure} goes
above the whole grid. Selecting a cell again returns to what is plotted there.

\ex{subplots.sci}
\fig{subplots.pdf}{0.95}{Four views of a damped oscillator, from
\texttt{subplots.sci}.}

A grid is saved and shown as one figure. In the IDE each cell is an interactive
chart of its own, and the plot pane's \textbf{Copy} and \textbf{Save\ldots} take
the whole grid.

\section{The figure as a record}

\code{plot.gcf()} returns the current figure as a record: its size, title,
labels and a list of series, each a record of its data and style. A figure kept
this way can be saved later with \code{plot.save(path, figure = f)}, after
other figures have been drawn.

\begin{axon}
import plot
plot.figure(title = "Kept")
plot.line([1; 2; 3], [3; 1; 2], label = "data")
kept = plot.gcf()
println(kept.title, length(kept.series), kept.series[1].type, kept.series[1].label)
plot.figure(title = "Another")
plot.bar([1; 2])
plot.save("_kept.svg", figure = kept)
import io
println(io.delete_file("_kept.svg"))
\end{axon}
\begin{axonout}
Kept 1 line data
true
\end{axonout}

The plot module is written in Axon (\texttt{stdlib/plot.sci}) over the
\code{graphics} module, which draws lists of lines, polygons, rectangles,
markers and text to PNG, SVG or PDF. \code{graphics} can be used directly for
drawings that are not charts; the Language Reference describes its scenes.

% ===================================================================
\chapter{Writing a module}\label{ch:module}

A module is an ordinary \code{.sci} file whose functions other programs can
import. This chapter builds a small statistics module: the mean, variance and
standard deviation of a sample, and a least-squares straight line. The same
steps apply to any module of your own.

\section{The module file}

A module is named after its file, so the file \texttt{stats.sci} is the module
\code{stats}. Here is the whole of \texttt{stats.sci}:

\ex{stats.sci}

The file shows the parts every module has:

\begin{itemize}
\item \textbf{Documentation.} The \code{##} block that opens the file describes
  the module, and the \code{##} lines above each function describe that
  function. They are what \code{help} shows, so write them for the person
  calling the function.
\item \textbf{Exports.} Only functions marked \code{export} can be used by
  programs that import the module: \code{mean}, \code{variance}, \code{std},
  \code{linfit} and \code{predict}. The helper \code{sample} is private, and
  reading \code{stats.sample} from outside raises \code{ImportError}. Keeping
  helpers private lets you change them later without breaking anyone's code.
\item \textbf{Checking arguments.} \code{sample} rejects anything that is not a
  vector or range, and data that is too short, with
  \code{error(kind, message)}. The message starts with the function's full name, so an error
  says where it came from. Use the standard kinds (\code{TypeError},
  \code{ArgumentError}, \code{DimensionError}, \ldots) so that callers can
  catch them by kind.
\item \textbf{Keyword parameters} give options with sensible defaults:
  \code{variance(x)} is the sample variance, and
  \code{variance(x, population = true)} the population variance.
\item \textbf{Records as results.} \code{linfit} returns several values in one
  record, \code|{slope, intercept, r2}|, which reads clearly at the call site
  (\code{fit.slope}) and can be passed on whole, as \code{predict} shows.
\end{itemize}

A module may define names that also exist as built-in functions. Inside
\texttt{stats.sci}, \code{mean} is the module's own function; programs reach it
as \code{stats.mean}, and the built-in \code{mean} is unaffected.

\section{Where to put the module file}\label{sec:modwhere}

When a program runs \code{import stats}, Axon looks for \texttt{stats.sci}
in three places, in this order, and uses the first it finds:

\begin{enumerate}
\item \textbf{In the same folder as your program.} This is the
  simplest choice, and the one this chapter uses: \texttt{stats.sci} sits next
  to \texttt{use\_stats.sci}. Nothing needs setting up, and the module stays
  with the programs when you copy the folder.

\begin{shell}
analysis/
  stats.sci           the module
  use_stats.sci       import stats
  experiment2.sci     import stats
\end{shell}

  The folder that counts is the one holding the \emph{importing file}, not
  the folder you start \texttt{axon} from, so a program finds its modules
  wherever it is run from. In the IDE, running a new program first asks
  you where to save it. Save it in the same folder as \texttt{stats.sci}, or
  \code{import stats} will not find the module.

  A subfolder groups related modules under one name. The module
  \texttt{models/kinetics.sci} is imported as \code{models.kinetics}, and
  \code{import models.kinetics as kin} gives it a short name.

\item \textbf{A shared folder listed in \texttt{package.toml}.} To use one copy
  of a module from programs in different folders, keep it in a folder of its
  own and name that folder in a file called \texttt{package.toml}, placed next
  to the programs (or in any folder above them; Axon uses the nearest one it
  finds going up). Paths are relative to the \texttt{package.toml} file.

\begin{shell}
shared/
  mystats/
    mystats.sci       import mystats
    extra.sci         import mystats.extra
project/
  package.toml
  main.sci
\end{shell}

  with \texttt{project/package.toml} containing

\begin{shell}
[dependencies]
mystats = "../shared/mystats"
\end{shell}

  \code{import mystats} then loads \texttt{shared/mystats/mystats.sci}, the
  folder's main module, and \code{import mystats.extra} loads
  \texttt{shared/mystats/extra.sci}. Fix a bug there once and every project
  that lists the folder gets the fix.

\item \textbf{The standard library.} Last, Axon looks in the \texttt{stdlib}
  folder next to the interpreter, where \code{plot}, \code{linalg} and
  \code{ode} live. You could add your own module there, but a rebuild or a new
  version of Axon replaces that folder, so it is better kept for modules that
  ship with Axon.
\end{enumerate}

Because the program's own folder is searched first, a module there with the
same name as a standard one (a \texttt{plot.sci} of your own, say) hides the
standard module for that program. Give your modules names that the standard
library does not use. If no place has the module, \code{import} raises
\code{ImportError}. Section~\ref{sec:modules} and the Language Reference's
``Modules'' section give the full rules.

\section{Using the module}

\ex{use_stats.sci}
\begin{axonout}
mean: 1.713
variance: 0.00412 population: 0.00353
std: 0.0642
slope: 1.9976 intercept: 0.0357 r2: 0.99884
\end{axonout}

\code{import stats} makes the module available as \code{stats}, so its
functions are called as \code{stats.mean} and so on.
\code{from stats import linfit, predict} binds chosen functions directly,
without the prefix. The
figure the program saves is Figure~\ref{fig:stats_fit.pdf}.

\fig{stats_fit.pdf}{0.6}{A least-squares line from \code{stats.linfit}, drawn
by \texttt{use\_stats.sci}.}

\section{Documentation and help}

Because the module has \code{##} comments, \code{help} works for it as for the
standard library, in the REPL and with \key{F1} in the IDE:

\begin{axonrepl}
> import stats
> help stats
module stats
  Basic statistics: the mean, variance and standard deviation of a sample,
  and a least-squares straight line through data points.
  exports: mean, variance, std, linfit, predict
> help stats.linfit
linfit(x, y)
  The least-squares straight line y = slope * x + intercept through the
  points (x[k], y[k]). Returns a record with slope, intercept and r2, the
  fraction of the variance of y that the line explains.
\end{axonrepl}

\section{Testing a module}

Check a new module against values you know. Here the built-in \code{mean}, a
variance worked by hand and an exact line give the answers to compare with;
\code{==} and a tolerance do the comparing.

\begin{axon}
import stats

x = [2.0; 4.0; 4.0; 4.0; 5.0; 5.0; 7.0; 9.0]
println(stats.mean(x) == mean(x))
println(stats.variance(x, population = true), stats.std(x, population = true))

# Points exactly on y = 3x - 1 give that line, with r2 = 1.
fit = stats.linfit(1:5, 3 * [1:5] - 1)
println(abs(fit.slope - 3) < 1e-12, abs(fit.intercept + 1) < 1e-12, fit.r2)

# Bad input raises errors you can catch.
try
  stats.variance({1, 2})
catch err
  println(err.kind, "-", err.message)
end
\end{axon}
\begin{axonout}
true
4.0 2.0
true true 1.0
TypeError - stats.variance: data must be a vector, got list
\end{axonout}

\section{Growing a module}

\begin{itemize}
\item \textbf{More files.} A folder of modules is a namespace:
  \texttt{mystats/stats.sci} and \texttt{mystats/tests.sci} are imported as
  \code{mystats.stats} and \code{mystats.tests}
  (\code{import mystats.stats as stats} keeps calls short).
\item \textbf{Sharing between projects.} Move the module to a folder of its
  own and list it in \texttt{package.toml}
  (Section~\ref{sec:modwhere}), rather than copying the file into each
  project.
\item \textbf{Speed.} Module functions are compiled by the JIT like any other
  (Chapter~\ref{ch:perf}).
\end{itemize}

% ===================================================================
\chapter{Performance}\label{ch:perf}

\section{The JIT compiler}

Axon compiles functions to native code at their first call, with a
just-in-time compiler built on the MIR library. Compilation changes only speed:
a compiled function gives the same results, prints the same output and raises
the same errors as when it is interpreted. A function that cannot be compiled
with that guarantee simply runs interpreted.

Loops and recursive functions typically run hundreds of times faster compiled.
These are Axon's own benchmarks, the best of five runs on a desktop PC, with
identical results both ways:

\begin{center}
\begin{tabular}{@{}lrrr@{}}
\toprule
Program & Interpreted & Compiled & Faster \\
\midrule
Recursive \texttt{fib(22)} & 132 ms & 222 \textmu s & 595\texttimes \\
A sum loop, 100\,000 terms & 89 ms & 106 \textmu s & 841\texttimes \\
Sieve of Eratosthenes to 100\,000 & 375 ms & 724 \textmu s & 517\texttimes \\
Mandelbrot set, 60 \texttimes{} 60 points & 190 ms & 212 \textmu s &
  896\texttimes \\
60 \texttimes{} 60 matrix product, as loops & 462 ms & 411 \textmu s &
  1123\texttimes \\
\texttt{ode.solve}, RK45, Lotka--Volterra & 22.6 ms & 5.63 ms & 4\texttimes \\
\bottomrule
\end{tabular}
\end{center}

\paragraph{What compiles.} Functions (named, anonymous and nested) whose bodies
use numbers, \code{bool}s and real or integer arrays: arithmetic and comparisons,
\code{if}, \code{while} and \code{for} loops, element reads and writes with
integer indices, calls to other functions, and the scalar built-ins (the
elementary functions, \code{abs}, \code{floor}, \code{min}, \code{max},
\code{mod}, \code{length}, \code{size}, \code{zeros}, \code{sum},
\code{println} and others). A parameter annotated \code{int}, \code{real},
\code{bool}, \code|vector{real}|, \code|vector{int}| or \code|matrix{real}| has
that type; an unannotated one takes the type of the argument it is called with,
and the function is compiled for each combination of argument types it meets.

\paragraph{What stays interpreted.} Functions that use strings (other than in
\code{print} and \code{println}), lists, records, dictionaries, complex numbers,
\code{try}, \code{global}, or indexing with ranges, colons or masks, and
functions whose variables change type. They still run correctly, at the
interpreter's speed.

\paragraph{Seeing what happened.} In the IDE, the Output pane lists after each
run the functions that were compiled and, for the others, why they were not.
The \textbf{JIT} button switches compilation off and on, to compare. From a
terminal, \code{axon --jit=off file.sci} runs without compiling.

\paragraph{Writing fast code.} Built-in array operations (\code{A * B},
\code{A \ b}, \code{sum}, elementwise arithmetic) are already native and fast.
Write loops where an algorithm needs them, keep each variable to one type, and
put the work in functions, since the JIT compiles functions, not top-level
code.

\section{Timing code}\label{sec:bench}

\ex{timing.sci}

\code{time.bench} reports the best time per call over several runs, which is
the measurement least disturbed by other activity on the machine. The program
prints, for example, \texttt{basel(1e5): 104 \textmu s per call} on the PC
used for the benchmarks above; run it with the JIT off and the loop takes
several hundred times longer.

% ===================================================================
\chapter{Calling C}\label{ch:c}

An \code{extern} block declares functions of a native library, which can then be
called like Axon functions. Arguments are converted to the declared C types
(\code{int32}, \code{float64}, \code{cstring} and so on) before the call, and a
value that does not fit is a \code{TypeError}. \code{name "cos"} calls a C
function under another Axon name.

\ex{extern_libc.sci}
\begin{axonout}
0.877582561890373 0.877582561890373
4
42
\end{axonout}

The library name is looked up in the nearest \texttt{package.toml}, which maps it
to a file on each platform, and otherwise on the system's library path:

\lstinputlisting[style=shellstyle,title={\small\ttfamily package.toml}]
  {examples/package.toml}

\paragraph{Arrays and pointers.} \code{ptr float64} passes a real array's buffer
directly, so C code can read and write it, and \code{ref T} passes a variable
by reference. \code{callback(...)} parameters accept Axon functions, which C
code can call back during the call. The standard library's \code{linalg} module
is written this way over OpenBLAS's LAPACKE interface; its source
(\texttt{stdlib/linalg.sci}) is a good model. The Language Reference's section
on foreign functions has the full rules.

% ===================================================================
\appendix
\chapter{Built-in functions}\label{app:builtins}

These functions and constants are available in every program without an
\code{import}. At the REPL, \code{help name} describes any of them, and in the
IDE \key{F1} does.

\section*{Arrays}
\begin{longtable}{@{}p{0.36\linewidth}p{0.58\linewidth}@{}}
\toprule
Function & Result \\
\midrule
\endhead
\code{zeros(m: int, n: int)}, \code{zeros(n: int)} & An m\texttimes n real matrix, or a vector,
  of zeros \\
\code{ones(m: int, n: int)}, \code{ones(n: int)} & The same, filled with ones \\
\code{eye(n: int)}, \code{eye(m: int, n: int)} & The identity matrix \\
\code{rand()}, \code{rand(n: int)}, \code{rand(m: int, n: int)} & Uniform random reals in
  $[0, 1)$ \\
\code{linspace(a: real, b: real, n: int)} & n equally spaced points from a to b \\
\code{range(n)}, \code{range(a, b)}, \code{range(a, b, h)} & \code{1:n}, or the
  range \code{a:h:b} in reading order (step 1 if left out) \\
\code{size(A)}, \code{size(A, d: int)} & The row and column counts as two values,
  or the count along dimension d \\
\code{length(x)} & The number of elements of an array, list, string or
  dictionary \\
\code{reshape(A, m: int, n: int)}, \code{reshape(A, n: int)} & The elements in column-major
  order, as an m\texttimes n matrix or a vector \\
\code{eltype(A)} & The element type: \code{"bool"}, \code{"int"},
  \code{"real"} or \code{"complex"} \\
\code{sum(A)}, \code{sum(A, d: int)} & The sum of all elements, or along dimension
  d \\
\code{prod(A)}, \code{prod(A, d: int)} & The product \\
\code{mean(A)}, \code{mean(A, d: int)} & The mean \\
\code{min(A)}, \code{max(A)} & The smallest or largest element \\
\code{min(a, b)}, \code{max(a, b)} & Elementwise, with broadcasting \\
\code{all(A)}, \code{any(A)} & Whether every, or some, element of a
  \code{bool} array is \code{true} \\
\code{dot(u, v)} & The inner product, with \code{u} conjugated \\
\code{tr(A)} & The transpose, the same as \code{A'} \\
\code{trace(A)} & The sum of the diagonal of a square matrix \\
\code{adjoint(A)} & The conjugate transpose \\
\code{sort(x)}, \code{sort(x, "descend")} & The elements in order, as a new
  vector, row or list \\
\code{sortperm(x)}, \code{sortperm(x, "descend")} & The positions \code{p}
  with \code{x[p]} sorted \\
\bottomrule
\end{longtable}

\section*{Sparse matrices}
\begin{longtable}{@{}p{0.36\linewidth}p{0.58\linewidth}@{}}
\toprule
Function & Result \\
\midrule
\endhead
\code{sparse(i, j, v, m: int, n: int)} & An m\texttimes n sparse matrix with entries
  \code{v[k]} at \code{(i[k], j[k])} \\
\code{sparse(A)} & A dense matrix as a sparse one \\
\code{speye(n: int)}, \code{speye(m: int, n: int)} & The sparse identity \\
\code{nnz(S)} & The number of stored entries \\
\code{full(S)} & A sparse matrix as a dense one \\
\code{splu(S)} & \code{L, U} with \code{S = L*U} \\
\bottomrule
\end{longtable}

\section*{Elementary functions}

\code{sqrt}, \code{exp}, \code{log} (natural), \code{sin}, \code{cos},
\code{tan}, \code{asin}, \code{acos}, \code{atan}, \code{sinh}, \code{cosh},
\code{tanh}, \code{asinh}, \code{acosh} and \code{atanh} take a number or an
array and apply elementwise. A real argument outside the domain gives NaN; a
complex one gives the principal value. \code{abs(x)} is the absolute value, or
the modulus of a complex number, and \code{factorial(n)} is $n!$ for an
\code{int} $n \ge 0$.

\section*{Rounding and integers}
\begin{longtable}{@{}p{0.36\linewidth}p{0.58\linewidth}@{}}
\toprule
Function & Result \\
\midrule
\endhead
\code{int(x)} & Converts to \code{int}, truncating toward zero \\
\code{floor(x)}, \code{ceil(x)} & The integer below or above, as an \code{int} \\
\code{round(x)} & The nearest integer, halves away from zero \\
\code{div(a, b)} & Floored division \\
\code{mod(a, b)} & The remainder of \code{div}, with the sign of \code{b} \\
\bottomrule
\end{longtable}

\section*{Complex numbers}
\begin{longtable}{@{}p{0.36\linewidth}p{0.58\linewidth}@{}}
\toprule
Function & Result \\
\midrule
\endhead
\code{complex(re: real, im: real)}, \code{complex(x)} & A complex number \\
\code{real(z)}, \code{imag(z)} & The real or imaginary part \\
\code{conj(z)} & The complex conjugate \\
\code{arg(z)} & The phase angle in radians \\
\bottomrule
\end{longtable}

\section*{Output and text}
\begin{longtable}{@{}p{0.36\linewidth}p{0.58\linewidth}@{}}
\toprule
Function & Result \\
\midrule
\endhead
\code{print(a, b, ...)} & Writes the values separated by spaces \\
\code{println(a, b, ...)} & The same, then ends the line \\
\code{format(x, d: int)} & \code{x} as text with exactly \code{d} decimals \\
\code{string(x)} & The display form of any value, as a string \\
\code{split(s: string)}, \code{split(s: string, sep: string)} & The words of s, or the pieces between
  each sep, as a list of strings \\
\code{join(parts: list, sep: string)} & A list of strings joined, with sep between them \\
\code{trim(s: string)}, \code{ltrim(s: string)}, \code{rtrim(s: string)} & s without spaces, tabs and
  line breaks at both ends, the start or the end \\
\code{upper(s: string)}, \code{lower(s: string)} & s in capital or small letters \\
\code{contains(s: string, p: string)}, \code{startswith(s: string, p: string)}, \code{endswith(s: string, p: string)} & Whether
  p occurs in s, begins it or ends it \\
\code{find(s: string, p: string)}, \code{find(s, p, start: int)} & The position of the first p in s,
  or 0 \\
\code{replace(s: string, old: string, new: string)} & s with every old replaced by new \\
\code{repeat(s: string, n: int)} & s repeated n times \\
\code{parse_function(params: string, expr: string)} & A function made from text
  (Section~\ref{sec:parsefn}) \\
\code{int(s: string)}, \code{real(s: string)} & The number written in the string s \\
\bottomrule
\end{longtable}

\section*{Lists and dictionaries}
\begin{longtable}{@{}p{0.36\linewidth}p{0.58\linewidth}@{}}
\toprule
Function & Result \\
\midrule
\endhead
\code{append(l, x)} & A new list or vector with \code{x} added \\
\code{dict(key => value, ...)} & A dictionary \\
\code{get(d: dict, key, default)} & \code{d[key]}, or \code{default} when missing \\
\code{haskey(d: dict, key)} & Whether the key is present \\
\code{keys(d: dict)} & The keys, as a list in insertion order \\
\bottomrule
\end{longtable}

\section*{Other functions}
\begin{longtable}{@{}p{0.36\linewidth}p{0.58\linewidth}@{}}
\toprule
Function & Result \\
\midrule
\endhead
\code{typeof(x)} & The type of a value, as a string \\
\code{isequal(a, b)} & Whether two values are equal as whole values \\
\code{isa(x, T)} & Whether x is exactly of type T: a type name such as
  \code|"int"| or \code|"vector{real}"|, or a declared record type \\
\code{error(message: string)}, \code{error(kind: string, message: string)} & Raises an error \\
\code{keep(f: function)} & Keeps a callback valid after the \code{extern} call that
  received it \\
\code{atexit(f: function)} & Calls \code{f()} when the program ends normally, the last
  registered first \\
\bottomrule
\end{longtable}

\section*{Constants}
\begin{longtable}{@{}p{0.36\linewidth}p{0.58\linewidth}@{}}
\toprule
Name & Value \\
\midrule
\endhead
\code{pi} & $\pi$, 3.141592653589793 \\
\code{e} & The base of the natural logarithm, 2.718281828459045 \\
\bottomrule
\end{longtable}

Like the functions, these are ordinary names: a program may assign \code{e} a
value of its own, which hides the constant from then on.

% ===================================================================
\chapter{The example programs}\label{app:examples}

These programs come with Axon and are listed on the IDE's Welcome tab, where each
opens as your own copy to change and run.

\newcommand{\shipped}[1]{\lstinputlisting[style=axonstyle,
  title={\small\ttfamily\detokenize{#1}}]{../../examples/#1}}

\shipped{01_basics.sci}
\shipped{02_collections.sci}
\shipped{03_linear_algebra.sci}
\shipped{04_lotka_volterra.sci}
\shipped{05_stiff_and_events.sci}
\shipped{06_sparse.sci}
\shipped{07_surface.sci}
\shipped{08_timing.sci}
\shipped{09_jit.sci}
\shipped{10_random.sci}
\shipped{11_bar_charts.sci}
\shipped{12_scatter.sci}
\shipped{13_line_plots.sci}
\shipped{14_lorenz.sci}
\shipped{15_subplots.sci}
\shipped{16_four_views.sci}

\end{document}
